\documentclass[10pt]{article}
% Two builds from this one file (see paper/README.md):
%   latexmk -pdf main.tex                     -> main.pdf, preprint with the author's name (for the review)
%   the anonymous build (defines \anonymousbuild) -> main_anonymous.pdf, TMLR double-blind submission format:
%   "Anonymous authors", no repository URL, a concise future-work paragraph without a timeline.
\newif\ifanonymous
\ifdefined\anonymousbuild\anonymoustrue\fi
\ifanonymous\usepackage{tmlr}\else\usepackage[preprint]{tmlr}\fi

\usepackage[utf8]{inputenc}
\usepackage[T1]{fontenc}
\usepackage{amsmath,amssymb}
\usepackage{booktabs}
\usepackage{graphicx}
\usepackage[export]{adjustbox}  % max width: figures are drawn at their printed size and only shrunk if needed
\usepackage{xspace}
\usepackage{multirow}
\usepackage{array}
\usepackage{xcolor}
\usepackage{hyperref}
\usepackage{url}
\hypersetup{colorlinks=true, linkcolor=blue!50!black, citecolor=blue!50!black, urlcolor=blue!50!black}

% Every result number is a macro generated from runs/, results/ and (post-freeze follow-up) explore/runs/ by
% paper/compute_numbers.py (see NUMBERS.md).
% AUTO-GENERATED by paper/compute_numbers.py from runs/ and results/. Do not edit by hand.
\newcommand{\framesPerTaskFlat}{150k\xspace} % frames per task, CNN/SNN fetch3 and fetch5 protocol
\newcommand{\framesPerTaskHybrid}{400k\xspace} % frames per task, full hybrid protocol
\newcommand{\framesPerTaskFiveFair}{450k\xspace} % frames per task, fetch5 SNN fair-budget runs
\newcommand{\framesDoorKey}{400k\xspace} % training frames, DoorKey-6x6 (RQ1)
\newcommand{\lrFlat}{0.001\xspace} % Adam learning rate, CNN/SNN trunks
\newcommand{\lrHybrid}{0.0003\xspace} % Adam learning rate, full hybrid
\newcommand{\numEnvs}{16\xspace} % parallel environments
\newcommand{\numSteps}{128\xspace} % rollout length per env per PPO update
\newcommand{\updateEpochs}{4\xspace} % PPO epochs per update
\newcommand{\minibatch}{256\xspace} % PPO minibatch size
\newcommand{\targetKl}{0.03\xspace} % PPO early-stop KL threshold
\newcommand{\clipCoef}{0.2\xspace} % PPO clip coefficient
\newcommand{\gammaVal}{0.99\xspace} % discount
\newcommand{\gaeLambda}{0.95\xspace} % GAE lambda
\newcommand{\snnT}{4\xspace} % SNN timesteps per environment step
\newcommand{\featDim}{128\xspace} % encoder feature dimension
\newcommand{\hybridWindow}{4\xspace} % Transformer context window (frames), hybrid
\newcommand{\paramsCnnShared}{194,312\xspace} % parameters, shared CNN agent (3 task heads)
\newcommand{\paramsCnnPerTask}{159,864\xspace} % parameters, one isolated CNN network
\newcommand{\paramsHybridShared}{327,800\xspace} % parameters, shared full-hybrid agent
\newcommand{\bufferPerTask}{5,000\xspace} % replay states stored per task (default)
\newcommand{\sleepPeriod}{25\xspace} % PPO updates between sleep phases (Sleep default)
\newcommand{\sleepSteps}{700\xspace} % gradient steps per sleep phase (Sleep default)
\newcommand{\evalEpisodesCl}{100\xspace} % evaluation episodes per task (continual)
\newcommand{\evalEpisodesSingle}{200\xspace} % evaluation episodes (single-task runs)
\newcommand{\homeoTarget}{0.15\xspace} % homeostasis target firing rate
\newcommand{\stdpAlpha}{0.03\xspace} % local-term weight alpha, stabilised STDP and random control
\newcommand{\stdpAlphaLarge}{0.3\xspace} % larger alpha tested (context arm tfs0.3)
\newcommand{\homeoLr}{0.01\xspace} % homeostasis bias learning rate (HybridSTDP default)
\newcommand{\convChannels}{16, 32 and 64\xspace} % conv channels (2x2 kernels) of both encoders
\newcommand{\lifBeta}{0.9\xspace} % initial LIF leak (learnable)
\newcommand{\surrogateSlope}{25\xspace} % surrogate-gradient slope
\newcommand{\solveThreshold}{0.5\xspace} % solved = final eval return above this
\newcommand{\memWindow}{12\xspace} % frame-stack length / context window in the memory probes
\newcommand{\sSevenFrames}{1M\xspace} % MemoryS7 pilot training frames
\newcommand{\pjMac}{4.6\xspace} % energy per 32-bit MAC used in the proxy (literature value, 45 nm)
\newcommand{\pjAc}{0.9\xspace} % energy per accumulate used in the proxy (literature value, 45 nm)
\newcommand{\bytesPerState}{188\xspace} % replay-buffer bytes per stored state, CNN/SNN trunks
\newcommand{\bytesPerStateHybrid}{632\xspace} % replay-buffer bytes per stored state, hybrid (4-frame window)
\newcommand{\rqThreeSeeds}{2\xspace} % seeds per RQ3 point
\newcommand{\rqThreeCnnDenseOps}{210k\xspace} % CNN event-driven ops/frame, no sparsity penalty
\newcommand{\rqThreeCnnDenseAcc}{0.999\xspace} % CNN BC test accuracy, no penalty
\newcommand{\rqThreeCnnDenseEnergy}{1653\xspace} % CNN dense energy proxy (nJ/frame), 4.6 pJ/MAC
\newcommand{\rqThreeCnnFloorOps}{33k\xspace} % lowest-ops CNN that still works (lam 10)
\newcommand{\rqThreeCnnFloorAcc}{0.998\xspace} % its accuracy
\newcommand{\rqThreeCnnCollapseOps}{6.9k\xspace} % CNN at lam 20: whole net silent
\newcommand{\rqThreeCnnCollapseAcc}{0.524\xspace} % its accuracy (= majority-action rate)
\newcommand{\rqThreeSnnTwoLowOps}{16.6k\xspace} % SNN T=2 at lam 0.03: ops/frame
\newcommand{\rqThreeSnnTwoLowAcc}{0.986\xspace} % SNN T=2 at lam 0.03: accuracy
\newcommand{\rqThreeSnnTwoLowRet}{0.970\xspace} % SNN T=2 at lam 0.03: closed-loop return
\newcommand{\rqThreeSnnTwoLowEnergy}{14.9\xspace} % SNN T=2 lam 0.03 energy proxy (nJ/frame), 0.9 pJ/SynOp
\newcommand{\rqThreeSnnFourLowOps}{16.2k\xspace} % SNN T=4 at lam 0.1: ops/frame
\newcommand{\rqThreeSnnFourLowAcc}{0.969\xspace} % SNN T=4 at lam 0.1: accuracy
\newcommand{\rqThreeCnnFloorEnergy}{126\xspace} % CNN lam 10 event-driven energy proxy (nJ/frame)
\newcommand{\rqThreeOpsRatio}{2.0\xspace} % CNN floor ops / SNN T=2 low-ops point
\newcommand{\msCnn}{16\xspace} % CPU ms per PPO minibatch step (256), CNN
\newcommand{\msSnn}{80\xspace} % CPU ms per PPO minibatch step, SNN T=4
\newcommand{\msHybrid}{93\xspace} % CPU ms per PPO minibatch step, SNN+Transformer
\newcommand{\snnCpuSlowdown}{5.2\xspace} % SNN / CNN CPU training cost per step
\newcommand{\rqOneBpN}{10\xspace} % seeds, arm bp
\newcommand{\rqOneBpSolved}{8/10\xspace} % solved (final eval return > 0.5), arm bp
\newcommand{\rqOneBpAuc}{0.416 $\pm$ 0.232\xspace} % AUC (mean train return) mean $\pm$ SD, arm bp
\newcommand{\rqOneBpEval}{0.766\xspace} % mean final eval return, arm bp
\newcommand{\rqOneHomeoN}{10\xspace} % seeds, arm homeo
\newcommand{\rqOneHomeoSolved}{10/10\xspace} % solved (final eval return > 0.5), arm homeo
\newcommand{\rqOneHomeoAuc}{0.616 $\pm$ 0.113\xspace} % AUC (mean train return) mean $\pm$ SD, arm homeo
\newcommand{\rqOneHomeoEval}{0.917\xspace} % mean final eval return, arm homeo
\newcommand{\rqOneStdpN}{10\xspace} % seeds, arm tfs0.03
\newcommand{\rqOneStdpSolved}{8/10\xspace} % solved (final eval return > 0.5), arm tfs0.03
\newcommand{\rqOneStdpAuc}{0.515 $\pm$ 0.300\xspace} % AUC (mean train return) mean $\pm$ SD, arm tfs0.03
\newcommand{\rqOneStdpEval}{0.765\xspace} % mean final eval return, arm tfs0.03
\newcommand{\rqOneRandN}{10\xspace} % seeds, arm rands0.03
\newcommand{\rqOneRandSolved}{8/10\xspace} % solved (final eval return > 0.5), arm rands0.03
\newcommand{\rqOneRandAuc}{0.482 $\pm$ 0.323\xspace} % AUC (mean train return) mean $\pm$ SD, arm rands0.03
\newcommand{\rqOneRandEval}{0.742\xspace} % mean final eval return, arm rands0.03
\newcommand{\rqOneDfaN}{10\xspace} % seeds, arm dfae
\newcommand{\rqOneDfaSolved}{2/10\xspace} % solved (final eval return > 0.5), arm dfae
\newcommand{\rqOneDfaAuc}{0.076 $\pm$ 0.101\xspace} % AUC (mean train return) mean $\pm$ SD, arm dfae
\newcommand{\rqOneDfaEval}{0.188\xspace} % mean final eval return, arm dfae
\newcommand{\rqOneDfaHomeoN}{10\xspace} % seeds, arm dfae_homeo
\newcommand{\rqOneDfaHomeoSolved}{7/10\xspace} % solved (final eval return > 0.5), arm dfae_homeo
\newcommand{\rqOneDfaHomeoAuc}{0.402 $\pm$ 0.307\xspace} % AUC (mean train return) mean $\pm$ SD, arm dfae_homeo
\newcommand{\rqOneDfaHomeoEval}{0.645\xspace} % mean final eval return, arm dfae_homeo
\newcommand{\rqOneFrozenN}{3\xspace} % seeds, arm frozen
\newcommand{\rqOneFrozenSolved}{2/3\xspace} % solved (final eval return > 0.5), arm frozen
\newcommand{\rqOneFrozenAuc}{0.441 $\pm$ 0.302\xspace} % AUC (mean train return) mean $\pm$ SD, arm frozen
\newcommand{\rqOneFrozenEval}{0.769\xspace} % mean final eval return, arm frozen
\newcommand{\rqOneStdpLargeN}{3\xspace} % seeds, arm tfs0.3
\newcommand{\rqOneStdpLargeSolved}{2/3\xspace} % solved (final eval return > 0.5), arm tfs0.3
\newcommand{\rqOneStdpLargeAuc}{0.397 $\pm$ 0.295\xspace} % AUC (mean train return) mean $\pm$ SD, arm tfs0.3
\newcommand{\rqOneStdpLargeEval}{0.670\xspace} % mean final eval return, arm tfs0.3
\newcommand{\rqOneHebN}{3\xspace} % seeds, arm hebs0.03
\newcommand{\rqOneHebSolved}{1/3\xspace} % solved (final eval return > 0.5), arm hebs0.03
\newcommand{\rqOneHebAuc}{0.188 $\pm$ 0.304\xspace} % AUC (mean train return) mean $\pm$ SD, arm hebs0.03
\newcommand{\rqOneHebEval}{0.333\xspace} % mean final eval return, arm hebs0.03
\newcommand{\rqOneContrAAucCi}{\ensuremath{-}0.101 [\ensuremath{-}0.323, +0.121]\xspace} % contrast (a) tfs0.03 vs homeo: AUC diff [95% Welch CI]
\newcommand{\rqOneContrAAucP}{0.338\xspace} % contrast (a) AUC Welch p (raw)
\newcommand{\rqOneContrAAucHolm}{0.724\xspace} % contrast (a) AUC Welch p, Holm over 4 contrasts
\newcommand{\rqOneContrAFisherP}{0.474\xspace} % contrast (a) solve-rate Fisher p (raw)
\newcommand{\rqOneContrAFisherHolm}{1.000\xspace} % contrast (a) Fisher p, Holm
\newcommand{\rqOneContrASolveCi}{\ensuremath{-}0.20 [\ensuremath{-}0.51, +0.11]\xspace} % contrast (a) solve-rate difference [95% Newcombe CI]
\newcommand{\rqOneContrBAucCi}{+0.033 [\ensuremath{-}0.260, +0.326]\xspace} % contrast (b) tfs0.03 vs rands0.03: AUC diff [95% Welch CI]
\newcommand{\rqOneContrBAucP}{0.818\xspace} % contrast (b) AUC Welch p (raw)
\newcommand{\rqOneContrBAucHolm}{0.818\xspace} % contrast (b) AUC Welch p, Holm over 4 contrasts
\newcommand{\rqOneContrBFisherP}{1.000\xspace} % contrast (b) solve-rate Fisher p (raw)
\newcommand{\rqOneContrBFisherHolm}{1.000\xspace} % contrast (b) Fisher p, Holm
\newcommand{\rqOneContrBSolveCi}{+0.00 [\ensuremath{-}0.34, +0.34]\xspace} % contrast (b) solve-rate difference [95% Newcombe CI]
\newcommand{\rqOneContrCAucCi}{+0.200 [+0.024, +0.377]\xspace} % contrast (c) homeo vs bp: AUC diff [95% Welch CI]
\newcommand{\rqOneContrCAucP}{0.029\xspace} % contrast (c) AUC Welch p (raw)
\newcommand{\rqOneContrCAucHolm}{0.117\xspace} % contrast (c) AUC Welch p, Holm over 4 contrasts
\newcommand{\rqOneContrCFisherP}{0.474\xspace} % contrast (c) solve-rate Fisher p (raw)
\newcommand{\rqOneContrCFisherHolm}{1.000\xspace} % contrast (c) Fisher p, Holm
\newcommand{\rqOneContrCSolveCi}{+0.20 [\ensuremath{-}0.11, +0.51]\xspace} % contrast (c) solve-rate difference [95% Newcombe CI]
\newcommand{\rqOneContrDAucCi}{\ensuremath{-}0.134 [\ensuremath{-}0.372, +0.104]\xspace} % contrast (d) rands0.03 vs homeo: AUC diff [95% Welch CI]
\newcommand{\rqOneContrDAucP}{0.241\xspace} % contrast (d) AUC Welch p (raw)
\newcommand{\rqOneContrDAucHolm}{0.724\xspace} % contrast (d) AUC Welch p, Holm over 4 contrasts
\newcommand{\rqOneContrDFisherP}{0.474\xspace} % contrast (d) solve-rate Fisher p (raw)
\newcommand{\rqOneContrDFisherHolm}{1.000\xspace} % contrast (d) Fisher p, Holm
\newcommand{\rqOneContrDSolveCi}{\ensuremath{-}0.20 [\ensuremath{-}0.51, +0.11]\xspace} % contrast (d) solve-rate difference [95% Newcombe CI]
\newcommand{\homeoEarlySeeds}{3 vs 5 seeds\xspace} % session-2 snapshot seed counts (homeo vs bp)
\newcommand{\homeoEarlyP}{0.088\xspace} % session-2 snapshot: homeo (seeds 1-3) vs bp (seeds 1-5) AUC Welch p
\newcommand{\homeoEarlySolved}{3/3\xspace} % session-2 snapshot: homeo solved, seeds 1-3
\newcommand{\vanillaStdpAucCi}{\ensuremath{-}0.502 [\ensuremath{-}0.834, \ensuremath{-}0.170]\xspace} % DoorKey-5x5: vanilla STDP (alpha 1) minus backprop, AUC [95% CI]
\newcommand{\vanillaStdpP}{0.015\xspace} % its Welch p
\newcommand{\vanillaStdpSolved}{0/3\xspace} % vanilla STDP solved (5x5)
\newcommand{\vanillaBpSolved}{3/3\xspace} % backprop solved (5x5)
\newcommand{\vanillaStdpRate}{0.39\xspace} % vanilla STDP: encoder activity (spike rate) over the last 20 updates
\newcommand{\bpRate}{0.20\xspace} % backprop: encoder activity over the last 20 updates
\newcommand{\dfaSolveCi}{+0.50 [+0.07, +0.74]\xspace} % DFA+homeo minus DFA solve rate [95% Newcombe CI]
\newcommand{\dfaFisherP}{0.070\xspace} % Fisher p (raw)
\newcommand{\dfaFisherHolm}{0.070\xspace} % Fisher p, Holm over the 2 planned tests
\newcommand{\dfaAucCi}{+0.326 [+0.101, +0.551]\xspace} % DFA+homeo minus DFA AUC [95% Welch CI]
\newcommand{\dfaAucP}{0.009\xspace} % AUC Welch p (raw)
\newcommand{\dfaAucHolm}{0.017\xspace} % AUC Welch p, Holm
\newcommand{\dfaAucRatio}{5.3\xspace} % AUC ratio DFA+homeo / DFA
\newcommand{\dfaEarlySolvedHomeo}{3/3\xspace} % session-2 snapshot: DFA+homeo solved (seeds 1-3)
\newcommand{\dfaEarlySolvedDfa}{1/5\xspace} % session-2 snapshot: DFA solved (seeds 1-5)
\newcommand{\dfaEarlyAucP}{<0.001\xspace} % session-2 snapshot AUC Welch p
\newcommand{\dfaContinualAcc}{0.450 $\pm$ 0.026\xspace} % SNN fetch3 continual ACC, DFA+homeo encoder + sleep
\newcommand{\accCnnNaive}{0.555 $\pm$ 0.013\xspace} % Cnn fetch3 naive: ACC mean $\pm$ SD
\newcommand{\forgetCnnNaive}{0.631\xspace} % Cnn fetch3 naive: mean FORGET
\newcommand{\nCnnNaive}{3\xspace} % Cnn fetch3 naive: seeds
\newcommand{\accCnnEwc}{0.674 $\pm$ 0.030\xspace} % Cnn fetch3 ewc_lam100: ACC mean $\pm$ SD
\newcommand{\forgetCnnEwc}{\ensuremath{-}0.002\xspace} % Cnn fetch3 ewc_lam100: mean FORGET
\newcommand{\nCnnEwc}{3\xspace} % Cnn fetch3 ewc_lam100: seeds
\newcommand{\accCnnReplay}{0.970 $\pm$ 0.009\xspace} % Cnn fetch3 replay: ACC mean $\pm$ SD
\newcommand{\forgetCnnReplay}{0.002\xspace} % Cnn fetch3 replay: mean FORGET
\newcommand{\nCnnReplay}{3\xspace} % Cnn fetch3 replay: seeds
\newcommand{\accCnnSleep}{0.978 $\pm$ 0.001\xspace} % Cnn fetch3 sleep: ACC mean $\pm$ SD
\newcommand{\forgetCnnSleep}{0.001\xspace} % Cnn fetch3 sleep: mean FORGET
\newcommand{\nCnnSleep}{3\xspace} % Cnn fetch3 sleep: seeds
\newcommand{\accCnnIso}{0.979 $\pm$ 0.001\xspace} % Cnn fetch3 isolation: ACC mean $\pm$ SD
\newcommand{\forgetCnnIso}{\ensuremath{-}0.000\xspace} % Cnn fetch3 isolation: mean FORGET
\newcommand{\nCnnIso}{3\xspace} % Cnn fetch3 isolation: seeds
\newcommand{\sleepVsReplayCnnCi}{+0.008 [\ensuremath{-}0.015, +0.030]\xspace} % Cnn: sleep minus replay ACC [95% Welch CI] (unmatched budget)
\newcommand{\sleepVsReplayCnnP}{0.284\xspace} % Cnn: its Welch p
\newcommand{\accSnnNaive}{0.550 $\pm$ 0.130\xspace} % Snn fetch3 naive: ACC mean $\pm$ SD
\newcommand{\forgetSnnNaive}{0.321\xspace} % Snn fetch3 naive: mean FORGET
\newcommand{\nSnnNaive}{3\xspace} % Snn fetch3 naive: seeds
\newcommand{\accSnnEwc}{0.716 $\pm$ 0.199\xspace} % Snn fetch3 ewc_lam100: ACC mean $\pm$ SD
\newcommand{\forgetSnnEwc}{\ensuremath{-}0.013\xspace} % Snn fetch3 ewc_lam100: mean FORGET
\newcommand{\nSnnEwc}{3\xspace} % Snn fetch3 ewc_lam100: seeds
\newcommand{\accSnnReplay}{0.814 $\pm$ 0.162\xspace} % Snn fetch3 replay: ACC mean $\pm$ SD
\newcommand{\forgetSnnReplay}{0.023\xspace} % Snn fetch3 replay: mean FORGET
\newcommand{\nSnnReplay}{3\xspace} % Snn fetch3 replay: seeds
\newcommand{\accSnnSleep}{0.889 $\pm$ 0.003\xspace} % Snn fetch3 sleep: ACC mean $\pm$ SD
\newcommand{\forgetSnnSleep}{0.006\xspace} % Snn fetch3 sleep: mean FORGET
\newcommand{\nSnnSleep}{3\xspace} % Snn fetch3 sleep: seeds
\newcommand{\accSnnIso}{0.833 $\pm$ 0.073\xspace} % Snn fetch3 isolation: ACC mean $\pm$ SD
\newcommand{\forgetSnnIso}{0.014\xspace} % Snn fetch3 isolation: mean FORGET
\newcommand{\nSnnIso}{3\xspace} % Snn fetch3 isolation: seeds
\newcommand{\sleepVsReplaySnnCi}{+0.074 [\ensuremath{-}0.328, +0.476]\xspace} % Snn: sleep minus replay ACC [95% Welch CI] (unmatched budget)
\newcommand{\sleepVsReplaySnnP}{0.510\xspace} % Snn: its Welch p
\newcommand{\accHybNaive}{0.452 $\pm$ 0.016\xspace} % Hyb fetch3 naive: ACC mean $\pm$ SD
\newcommand{\forgetHybNaive}{0.536\xspace} % Hyb fetch3 naive: mean FORGET
\newcommand{\nHybNaive}{3\xspace} % Hyb fetch3 naive: seeds
\newcommand{\accHybReplay}{0.862 $\pm$ 0.038\xspace} % Hyb fetch3 replay: ACC mean $\pm$ SD
\newcommand{\forgetHybReplay}{0.019\xspace} % Hyb fetch3 replay: mean FORGET
\newcommand{\nHybReplay}{8\xspace} % Hyb fetch3 replay: seeds
\newcommand{\accHybSleep}{0.887 $\pm$ 0.045\xspace} % Hyb fetch3 sleep: ACC mean $\pm$ SD
\newcommand{\forgetHybSleep}{0.006\xspace} % Hyb fetch3 sleep: mean FORGET
\newcommand{\nHybSleep}{6\xspace} % Hyb fetch3 sleep: seeds
\newcommand{\accHybIso}{0.832 $\pm$ 0.037\xspace} % Hyb fetch3 isolation: ACC mean $\pm$ SD
\newcommand{\forgetHybIso}{0.006\xspace} % Hyb fetch3 isolation: mean FORGET
\newcommand{\nHybIso}{3\xspace} % Hyb fetch3 isolation: seeds
\newcommand{\accHybSleepMatched}{0.879 $\pm$ 0.043\xspace} % Hyb fetch3 sleep_matched: ACC mean $\pm$ SD
\newcommand{\forgetHybSleepMatched}{0.020\xspace} % Hyb fetch3 sleep_matched: mean FORGET
\newcommand{\nHybSleepMatched}{8\xspace} % Hyb fetch3 sleep_matched: seeds
\newcommand{\sleepVsReplayHybCi}{+0.024 [\ensuremath{-}0.026, +0.075]\xspace} % Hyb: sleep minus replay ACC [95% Welch CI] (unmatched budget)
\newcommand{\sleepVsReplayHybP}{0.309\xspace} % Hyb: its Welch p
\newcommand{\accCnnFiveNaive}{0.416 $\pm$ 0.098\xspace} % CnnFive fetch5 naive: ACC mean $\pm$ SD
\newcommand{\forgetCnnFiveNaive}{0.612\xspace} % CnnFive fetch5 naive: mean FORGET
\newcommand{\nCnnFiveNaive}{3\xspace} % CnnFive fetch5 naive: seeds
\newcommand{\accCnnFiveReplay}{0.926 $\pm$ 0.049\xspace} % CnnFive fetch5 replay: ACC mean $\pm$ SD
\newcommand{\forgetCnnFiveReplay}{0.006\xspace} % CnnFive fetch5 replay: mean FORGET
\newcommand{\nCnnFiveReplay}{3\xspace} % CnnFive fetch5 replay: seeds
\newcommand{\accCnnFiveSleep}{0.928 $\pm$ 0.065\xspace} % CnnFive fetch5 sleep: ACC mean $\pm$ SD
\newcommand{\forgetCnnFiveSleep}{0.003\xspace} % CnnFive fetch5 sleep: mean FORGET
\newcommand{\nCnnFiveSleep}{3\xspace} % CnnFive fetch5 sleep: seeds
\newcommand{\accCnnFiveIso}{0.976 $\pm$ 0.000\xspace} % CnnFive fetch5 isolation: ACC mean $\pm$ SD
\newcommand{\forgetCnnFiveIso}{\ensuremath{-}0.000\xspace} % CnnFive fetch5 isolation: mean FORGET
\newcommand{\nCnnFiveIso}{3\xspace} % CnnFive fetch5 isolation: seeds
\newcommand{\sleepVsReplayCnnFiveCi}{+0.002 [\ensuremath{-}0.133, +0.137]\xspace} % CnnFive: sleep minus replay ACC [95% Welch CI] (unmatched budget)
\newcommand{\sleepVsReplayCnnFiveP}{0.969\xspace} % CnnFive: its Welch p
\newcommand{\ewcAccTen}{0.816 $\pm$ 0.120\xspace} % CNN EWC lambda=10: ACC
\newcommand{\ewcForgetTen}{\ensuremath{-}0.005\xspace} % CNN EWC lambda=10: FORGET
\newcommand{\ewcLastTaskTen}{0.603\xspace} % CNN EWC lambda=10: return on the last task after training it
\newcommand{\ewcAccHundred}{0.674 $\pm$ 0.030\xspace} % CNN EWC lambda=100: ACC
\newcommand{\ewcForgetHundred}{\ensuremath{-}0.002\xspace} % CNN EWC lambda=100: FORGET
\newcommand{\ewcLastTaskHundred}{0.441\xspace} % CNN EWC lambda=100: return on the last task after training it
\newcommand{\ewcAccThousand}{0.545 $\pm$ 0.006\xspace} % CNN EWC lambda=1000: ACC
\newcommand{\ewcForgetThousand}{0.000\xspace} % CNN EWC lambda=1000: FORGET
\newcommand{\ewcLastTaskThousand}{0.280\xspace} % CNN EWC lambda=1000: return on the last task after training it
\newcommand{\naiveLastTaskCnn}{0.971\xspace} % CNN naive: return on the last task
\newcommand{\matchN}{8\xspace} % block 1 seeds per arm
\newcommand{\matchSamplesMin}{0.91M\xspace} % smallest per-run replay count
\newcommand{\matchSamplesMax}{1.48M\xspace} % largest per-run replay count
\newcommand{\unmatchedSleepSamples}{1.43M\xspace} % replay states consumed by an unmatched hybrid sleep run
\newcommand{\unmatchedExtraPct}{28\xspace} % % more replayed states for unmatched sleep vs replay (same seeds 1-6)
\newcommand{\matchAccSleep}{0.879 $\pm$ 0.043\xspace} % block 1 sleep_matched ACC
\newcommand{\matchAccReplay}{0.862 $\pm$ 0.038\xspace} % block 1 replay ACC
\newcommand{\matchAccCi}{+0.017 [\ensuremath{-}0.026, +0.060]\xspace} % block 1 sleep_matched minus replay ACC [95% Welch CI]
\newcommand{\matchAccP}{0.415\xspace} % block 1 ACC Welch p
\newcommand{\matchForgetSleep}{0.020 $\pm$ 0.011\xspace} % block 1 sleep_matched FORGET
\newcommand{\matchForgetReplay}{0.019 $\pm$ 0.017\xspace} % block 1 replay FORGET
\newcommand{\matchForgetCi}{+0.001 [\ensuremath{-}0.015, +0.017]\xspace} % block 1 sleep_matched minus replay FORGET [95% Welch CI]
\newcommand{\matchForgetP}{0.881\xspace} % block 1 FORGET Welch p
\newcommand{\unmatchedForgetSameSeeds}{0.006\xspace} % unmatched sleep FORGET, seeds 1-6
\newcommand{\matchedForgetSameSeeds}{0.022\xspace} % matched sleep FORGET, seeds 1-6
\newcommand{\matchedMinusUnmatchedAccCi}{\ensuremath{-}0.014 [\ensuremath{-}0.029, +0.002]\xspace} % matched minus unmatched sleep ACC, same seeds [95% paired-t CI] (descriptive)
\newcommand{\matchedMinusUnmatchedAccP}{0.069\xspace} % its paired t p
\newcommand{\hybSleepReplayNThreeCi}{+0.039 [\ensuremath{-}0.006, +0.084]\xspace} % hybrid unmatched sleep minus replay ACC, seeds 1-3 [95% CI]
\newcommand{\hybSleepReplayNThreeP}{0.072\xspace} % its Welch p (seeds 1-3)
\newcommand{\hybSleepReplayNThreeSeeds}{seeds 1--3\xspace} % seed set of the snapshot
\newcommand{\hybSleepReplayNSixCi}{+0.032 [\ensuremath{-}0.023, +0.087]\xspace} % hybrid unmatched sleep minus replay ACC, seeds 1-6 [95% CI]
\newcommand{\hybSleepReplayNSixP}{0.224\xspace} % its Welch p (seeds 1-6)
\newcommand{\hybSleepReplayNSixSeeds}{seeds 1--6\xspace} % seed set of the snapshot
\newcommand{\snnSleepVanillaStdpAcc}{0.163\xspace} % SNN fetch3 sleep_stdp: ACC
\newcommand{\snnSleepVanillaStdpN}{1\xspace} % SNN fetch3 sleep_stdp: seeds
\newcommand{\snnSleepStabStdpAcc}{0.881 $\pm$ 0.028\xspace} % SNN fetch3 sleep_stdps0.003: ACC
\newcommand{\snnSleepStabStdpN}{3\xspace} % SNN fetch3 sleep_stdps0.003: seeds
\newcommand{\snnSleepHomeoAcc}{0.845 $\pm$ 0.034\xspace} % SNN fetch3 sleep_homeo: ACC
\newcommand{\snnSleepHomeoN}{3\xspace} % SNN fetch3 sleep_homeo: seeds
\newcommand{\memCnnThreeIsoMb}{1.92\xspace} % CNN fetch3: isolation total MB
\newcommand{\memCnnThreeSleepMb}{3.60\xspace} % CNN fetch3: sleep (5000/task) total MB
\newcommand{\memCnnThreeIsoPerMb}{0.51\xspace} % CNN fetch3: isolation ACC per MB
\newcommand{\memCnnThreeSleepPerMb}{0.27\xspace} % CNN fetch3: sleep ACC per MB
\newcommand{\memCnnThreeIsoAcc}{0.979\xspace} % CNN fetch3: isolation ACC
\newcommand{\memCnnThreeSleepAcc}{0.978\xspace} % CNN fetch3: sleep ACC
\newcommand{\perParamCnnThreeIso}{0.20\xspace} % CNN fetch3: isolation ACC per 100k params
\newcommand{\perParamCnnThreeSleep}{0.50\xspace} % CNN fetch3: sleep ACC per 100k params
\newcommand{\memSnnThreeIsoMb}{1.92\xspace} % SNN fetch3: isolation total MB
\newcommand{\memSnnThreeSleepMb}{3.60\xspace} % SNN fetch3: sleep (5000/task) total MB
\newcommand{\memSnnThreeIsoPerMb}{0.43\xspace} % SNN fetch3: isolation ACC per MB
\newcommand{\memSnnThreeSleepPerMb}{0.25\xspace} % SNN fetch3: sleep ACC per MB
\newcommand{\memSnnThreeIsoAcc}{0.833\xspace} % SNN fetch3: isolation ACC
\newcommand{\memSnnThreeSleepAcc}{0.889\xspace} % SNN fetch3: sleep ACC
\newcommand{\perParamSnnThreeIso}{0.17\xspace} % SNN fetch3: isolation ACC per 100k params
\newcommand{\perParamSnnThreeSleep}{0.46\xspace} % SNN fetch3: sleep ACC per 100k params
\newcommand{\memHybThreeIsoMb}{3.52\xspace} % Hybrid fetch3: isolation total MB
\newcommand{\memHybThreeSleepMb}{10.79\xspace} % Hybrid fetch3: sleep (5000/task) total MB
\newcommand{\memHybThreeIsoPerMb}{0.24\xspace} % Hybrid fetch3: isolation ACC per MB
\newcommand{\memHybThreeSleepPerMb}{0.08\xspace} % Hybrid fetch3: sleep ACC per MB
\newcommand{\memHybThreeIsoAcc}{0.832\xspace} % Hybrid fetch3: isolation ACC
\newcommand{\memHybThreeSleepAcc}{0.887\xspace} % Hybrid fetch3: sleep ACC
\newcommand{\perParamHybThreeIso}{0.09\xspace} % Hybrid fetch3: isolation ACC per 100k params
\newcommand{\perParamHybThreeSleep}{0.27\xspace} % Hybrid fetch3: sleep ACC per 100k params
\newcommand{\memCnnFiveIsoMb}{3.20\xspace} % CNN fetch5: isolation total MB
\newcommand{\memCnnFiveSleepMb}{5.61\xspace} % CNN fetch5: sleep (5000/task) total MB
\newcommand{\memCnnFiveIsoPerMb}{0.31\xspace} % CNN fetch5: isolation ACC per MB
\newcommand{\memCnnFiveSleepPerMb}{0.17\xspace} % CNN fetch5: sleep ACC per MB
\newcommand{\memCnnFiveIsoAcc}{0.976\xspace} % CNN fetch5: isolation ACC
\newcommand{\memCnnFiveSleepAcc}{0.928\xspace} % CNN fetch5: sleep ACC
\newcommand{\perParamCnnFiveIso}{0.12\xspace} % CNN fetch5: isolation ACC per 100k params
\newcommand{\perParamCnnFiveSleep}{0.41\xspace} % CNN fetch5: sleep ACC per 100k params
\newcommand{\memSnnFiveFairIsoMb}{3.20\xspace} % SNN fetch5: isolation total MB
\newcommand{\memSnnFiveFairSleepMb}{5.62\xspace} % SNN fetch5: sleep (5000/task) total MB
\newcommand{\memSnnFiveFairIsoPerMb}{0.27\xspace} % SNN fetch5: isolation ACC per MB
\newcommand{\memSnnFiveFairSleepPerMb}{0.16\xspace} % SNN fetch5: sleep ACC per MB
\newcommand{\memSnnFiveFairIsoAcc}{0.871\xspace} % SNN fetch5: isolation ACC
\newcommand{\memSnnFiveFairSleepAcc}{0.915\xspace} % SNN fetch5: sleep ACC
\newcommand{\perParamSnnFiveFairIso}{0.11\xspace} % SNN fetch5: isolation ACC per 100k params
\newcommand{\perParamSnnFiveFairSleep}{0.40\xspace} % SNN fetch5: sleep ACC per 100k params
\newcommand{\perParamRatioRange}{2.5--3.3\xspace} % sleep / isolation ACC per parameter, range over CNN3, SNN3, hybrid3, CNN5
\newcommand{\perMbRatioRange}{1.8--2.9\xspace} % isolation / sleep(5000) ACC per MB, range over CNN3, SNN3, hybrid3, CNN5
\newcommand{\smallCnnTwoHundredAcc}{0.964 $\pm$ 0.009\xspace} % CNN fetch3 fetch3_sleep_buf200 @150k: ACC
\newcommand{\smallCnnTwoHundredMb}{0.89\xspace} % CNN fetch3 fetch3_sleep_buf200 @150k: total MB
\newcommand{\smallCnnTwoHundredPerMb}{1.08\xspace} % CNN fetch3 fetch3_sleep_buf200 @150k: ACC per MB
\newcommand{\smallSnnTwoHundredAcc}{0.838 $\pm$ 0.016\xspace} % SNN fetch3 fetch3_sleep_buf200 @150k: ACC
\newcommand{\smallSnnTwoHundredMb}{0.89\xspace} % SNN fetch3 fetch3_sleep_buf200 @150k: total MB
\newcommand{\smallSnnTwoHundredPerMb}{0.94\xspace} % SNN fetch3 fetch3_sleep_buf200 @150k: ACC per MB
\newcommand{\smallSnnTwoThousandAcc}{0.847 $\pm$ 0.083\xspace} % SNN fetch3 fetch3_sleep_buf2000 @150k: ACC
\newcommand{\smallSnnTwoThousandMb}{1.91\xspace} % SNN fetch3 fetch3_sleep_buf2000 @150k: total MB
\newcommand{\smallSnnTwoThousandPerMb}{0.44\xspace} % SNN fetch3 fetch3_sleep_buf2000 @150k: ACC per MB
\newcommand{\smallSnnReplayTwoThousandAcc}{0.828 $\pm$ 0.167\xspace} % SNN fetch3 fetch3_replay_buf2000 @150k: ACC
\newcommand{\smallSnnReplayTwoThousandMb}{1.91\xspace} % SNN fetch3 fetch3_replay_buf2000 @150k: total MB
\newcommand{\smallSnnReplayTwoThousandPerMb}{0.43\xspace} % SNN fetch3 fetch3_replay_buf2000 @150k: ACC per MB
\newcommand{\smallSnnReplayTwoHundredAcc}{0.792 $\pm$ 0.183\xspace} % SNN fetch3 fetch3_replay_buf200 @150k: ACC
\newcommand{\smallSnnReplayTwoHundredMb}{0.89\xspace} % SNN fetch3 fetch3_replay_buf200 @150k: total MB
\newcommand{\smallSnnReplayTwoHundredPerMb}{0.89\xspace} % SNN fetch3 fetch3_replay_buf200 @150k: ACC per MB
\newcommand{\snnSleepTwoHundredVsIsoCi}{+0.005 [\ensuremath{-}0.167, +0.177]\xspace} % SNN sleep_buf200 minus isolation ACC [95% CI]
\newcommand{\snnSleepTwoHundredVsIsoP}{0.921\xspace} % its Welch p
\newcommand{\snnSleepTwoThousandVsIsoCi}{+0.014 [\ensuremath{-}0.165, +0.193]\xspace} % SNN sleep_buf2000 minus isolation ACC [95% CI]
\newcommand{\snnSleepTwoThousandVsIsoP}{0.836\xspace} % its Welch p
\newcommand{\snnIsoAcc}{0.833 $\pm$ 0.073\xspace} % SNN fetch3 isolation ACC
\newcommand{\fiveSnnShortSleep}{0.778 $\pm$ 0.007\xspace} % fetch5 SNN (Short) sleep ACC
\newcommand{\fiveSnnShortIso}{0.475 $\pm$ 0.023\xspace} % fetch5 SNN (Short) isolation ACC
\newcommand{\fiveSnnShortCi}{+0.303 [+0.148, +0.457]\xspace} % fetch5 SNN (Short) sleep minus isolation [95% CI]
\newcommand{\fiveSnnShortP}{0.023\xspace} % fetch5 SNN (Short) Welch p
\newcommand{\fiveSnnShortN}{2\xspace} % fetch5 SNN (Short) seeds per arm
\newcommand{\fiveSnnFairSleep}{0.915 $\pm$ 0.014\xspace} % fetch5 SNN (Fair) sleep ACC
\newcommand{\fiveSnnFairIso}{0.871 $\pm$ 0.026\xspace} % fetch5 SNN (Fair) isolation ACC
\newcommand{\fiveSnnFairCi}{+0.044 [\ensuremath{-}0.009, +0.098]\xspace} % fetch5 SNN (Fair) sleep minus isolation [95% CI]
\newcommand{\fiveSnnFairP}{0.078\xspace} % fetch5 SNN (Fair) Welch p
\newcommand{\fiveSnnFairN}{3\xspace} % fetch5 SNN (Fair) seeds per arm
\newcommand{\learnThreshold}{0.8\xspace} % task counted as learned if R[k][k] >= this (pre-registered)
\newcommand{\fiveSnnFairLearnedSleep}{14/15\xspace} % fetch5 SNN fair: tasks learned to >= 0.8, sleep
\newcommand{\fiveSnnFairLearnedIso}{12/15\xspace} % fetch5 SNN fair: tasks learned to >= 0.8, isolation
\newcommand{\fiveSnnFairFwt}{0.306 $\pm$ 0.170\xspace} % fetch5 SNN fair: sleep FWT vs isolation curves
\newcommand{\fiveSnnFairFwtP}{0.089\xspace} % one-sample t p (FWT vs 0)
\newcommand{\fwtCnn}{\ensuremath{-}0.017 $\pm$ 0.141\xspace} % Cnn fetch3 sleep FWT (vs same-seed isolation curves)
\newcommand{\fwtCnnP}{0.854\xspace} % one-sample t p
\newcommand{\fwtSnn}{0.199 $\pm$ 0.079\xspace} % Snn fetch3 sleep FWT (vs same-seed isolation curves)
\newcommand{\fwtSnnP}{0.049\xspace} % one-sample t p
\newcommand{\fwtHyb}{0.423 $\pm$ 0.244\xspace} % Hyb fetch3 sleep FWT (vs same-seed isolation curves)
\newcommand{\fwtHybP}{0.095\xspace} % one-sample t p
\newcommand{\fwtCnnFive}{\ensuremath{-}0.525 $\pm$ 0.722\xspace} % CNN fetch5 sleep FWT
\newcommand{\sSevenCnnOne}{0.915\xspace} % MemoryS7 pilot mem_cnn_fs1: eval success (1 seed)
\newcommand{\sSevenCnnFour}{0.995\xspace} % MemoryS7 pilot mem_cnn_fs4: eval success (1 seed)
\newcommand{\sSevenCnnEight}{0.835\xspace} % MemoryS7 pilot mem_cnn_fs8: eval success (1 seed)
\newcommand{\sSevenSnnEight}{0.465\xspace} % MemoryS7 pilot mem_snn_fs8: eval success (1 seed)
\newcommand{\sSevenSnnTf}{0.970\xspace} % MemoryS7 pilot mem_snn_tf8: eval success (1 seed)
\newcommand{\sSevenLrTf}{0.0003\xspace} % S7 pilot SNN+Transformer lr
\newcommand{\sSevenLrOther}{0.001\xspace} % S7 pilot other arms lr
\newcommand{\probeElevenOne}{0.430\xspace} % probe_S11_cnn_fs1: eval success
\newcommand{\probeElevenOneLen}{9\xspace} % probe_S11_cnn_fs1: mean episode length
\newcommand{\probeElevenTwelve}{0.465\xspace} % probe_S11_cnn_fs12: eval success
\newcommand{\probeElevenTwelveLen}{10\xspace} % probe_S11_cnn_fs12: mean episode length
\newcommand{\probeThirteenOne}{0.485\xspace} % probe_S13_cnn_fs1: eval success
\newcommand{\probeThirteenOneLen}{11\xspace} % probe_S13_cnn_fs1: mean episode length
\newcommand{\probeThirteenTwelve}{0.460\xspace} % probe_S13_cnn_fs12: eval success
\newcommand{\probeThirteenTwelveLen}{12\xspace} % probe_S13_cnn_fs12: mean episode length
\newcommand{\probeElevenTwelveLong}{0.485\xspace} % probe_S11_cnn_fs12_6M: eval success
\newcommand{\probeElevenTwelveLongLen}{10\xspace} % probe_S11_cnn_fs12_6M: mean episode length
\newcommand{\probeLongRange}{0.48--0.51\xspace} % S11 fs12 6M: range of training return at 1..6M frames
\newcommand{\probeFrames}{2M\xspace} % probe training frames
\newcommand{\probeLongFrames}{6M\xspace} % exploratory long probe: training frames
\newcommand{\probeLr}{0.0003\xspace} % probe lr
\newcommand{\rqFourCnn}{0.978\xspace} % continual ACC per unit training compute, CNN
\newcommand{\rqFourSnn}{0.173\xspace} % same, SNN + sleep
\newcommand{\rqFourSnnHomeo}{0.164\xspace} % same, SNN + homeostasis
\newcommand{\rqFourHybrid}{0.056\xspace} % same, full hybrid
\newcommand{\rqFourStdp}{0.084\xspace} % same, stabilised STDP
\newcommand{\rqFourDfaHomeo}{0.090\xspace} % same, DFA + homeostasis
\newcommand{\rqFourRatio}{5.7\xspace} % CNN / best spiking variant, ACC per unit compute
\newcommand{\rqFourLocalOnly}{0/6\xspace} % local-only (STDP-trained) encoder solve rate
\newcommand{\demoNaiveRow}{0.126, 0.542, 0.956\xspace} % demo naive seed 1: final accuracy-matrix row
\newcommand{\demoSleepRow}{0.978, 0.979, 0.979\xspace} % demo sleep seed 1: final accuracy-matrix row
\newcommand{\seedRange}{3--10\xspace} % seeds per arm, range over the main arms
\newcommand{\computeCoreHours}{110\xspace} % total single-thread CPU hours of all frozen PPO runs (continual + single-task; excludes RQ3 behaviour cloning)
\newcommand{\computeRuns}{289\xspace} % number of frozen PPO runs counted in computeCoreHours

\xspaceaddexceptions{\%}

% Claim labels: literature claim / our interpretation / summaries of our results; and the evidential status of each
% result of ours: confirmatory / exploratory / descriptive, plus post-freeze for the follow-up obtained after the freeze.
\newcommand{\lit}{\textsc{[lit]}\xspace}
\newcommand{\ours}{\textsc{[ours]}\xspace}
\newcommand{\interp}{\textsc{[interp]}\xspace}
\newcommand{\conf}{\textsc{[confirmatory]}\xspace}
\newcommand{\expl}{\textsc{[exploratory]}\xspace}
\newcommand{\desc}{\textsc{[descriptive]}\xspace}
\newcommand{\postfz}{\textsc{[post-freeze]}\xspace}
\newcommand{\todo}[1]{\textcolor{red}{\textbf{TODO:} #1}}
\newcommand{\repourl}{\ifanonymous an anonymised repository provided as supplementary material\else
  \url{https://github.com/VT69/NeuroPlast}\fi}
\newcommand{\fig}[1]{\includegraphics[max width=\linewidth]{figures/#1.pdf}}

\title{Testing Brain-Inspired Mechanisms for\\ Task-Incremental Continual Reinforcement Learning:\\
A Controlled Ablation Study}

\author{\name Vaibhav Tiwari \\
      \addr \todo{affiliation and contact email}}

\begin{document}

\maketitle

\begin{abstract}
Brain-inspired mechanisms such as spiking neurons, local plasticity, homeostasis and sleep-like consolidation are
proposed as remedies for catastrophic forgetting in continual reinforcement learning (RL), usually several at
once. We ask which of them help when each is isolated and compared with a matched non-biological control. In
task-incremental MiniGrid sequences (the task identity is given) with small PPO agents, we test six mechanisms one at
a time, with tests prospectively specified in a version-controlled experiment log, \seedRange seeds per arm, 95\%
confidence intervals and matched budgets. A seventh, Transformer working memory, is inconclusive: no benchmark we
tried was shown to be valid (requiring memory and learnable within budget). No mechanism produced a statistically
reliable improvement in final accuracy or forgetting beyond its matched control; given the seed counts, these results
rule out large effects rather than show equivalence, and they do so only where the intervals are narrow. On
DoorKey-6x6, the tested stabilised reward-modulated STDP rule was not distinguishable from a random update of the same
size (AUC difference \rqOneContrBAucCi, an interval too wide to exclude large effects). Sleep-like consolidation
removed most of naive fine-tuning's forgetting but, at an exactly matched replay budget, was not distinguishable from
experience replay (final-accuracy difference \matchAccCi). Two narrower effects stand. Homeostasis sped up learning in
a spiking encoder trained with backprop-free encoder credit assignment (direct feedback alignment; the heads use exact
gradients), with an AUC difference of \dfaAucCi (Holm $p=\dfaAucHolm$), without making it reliable. With a small replay buffer, a
shared trunk with per-task heads and sleep reached similar mean accuracy to full-width isolation (one network per task;
\smallSnnTwoHundredMean vs \snnIsoMean; \smallSnnTwoHundredN seeds, wide CI) in less than half the total memory; at
the default buffer, isolation was more memory-efficient. A post-freeze follow-up qualifies this: on a 5-task CNN
sequence with \postN fresh seeds, isolation with quarter-width networks beat a shared trunk trained with LwF or with
replay at equal or lower memory (\postNarrowMean vs \postLwfMean and \postReplayMean), so the comparison depends on
isolation's network width. A spiking encoder works at fewer operations per frame than a sparsified CNN, but its
measured training cost in our CPU implementation is \snnCpuSlowdown{}$\times$ higher, and its energy advantage rests
on an estimated arithmetic energy proxy. Seven earlier positive findings shrank or disappeared under controls; we
report each with its cause.
\end{abstract}

% =====================================================================================================
\section{Introduction}
\label{sec:intro}

A continual reinforcement learning (RL) agent learns a sequence of tasks while its parameters keep changing. With
plain gradient-based fine-tuning, learning a new task overwrites what earlier tasks needed: catastrophic
forgetting. The problem is framed as a stability--plasticity trade-off \lit \citep{khetarpal2022towards,
kirkpatrick2017overcoming}, and the standard remedies are regularisation towards old parameters
\citep[EWC;][]{kirkpatrick2017overcoming}, replaying stored experience \citep[CLEAR;][]{rolnick2019experience} and
isolating parameters per task \citep[e.g. PackNet;][]{mallya2018packnet}.

A second family of remedies is inspired by the brain, and it is popular. Spiking neural networks (SNNs) promise
event-driven, low-energy computation \lit \citep{fang2023spikingjelly, zhou2026spikingformer}. Local learning rules
such as STDP and its reward-modulated three-factor variants are proposed as biologically plausible complements or
alternatives to backpropagation \lit \citep{lee2025brain, huo2025research}. Sleep-like offline replay with local
plasticity has been reported to recover forgotten tasks \lit \citep{tadros2022sleep}, and homeostatic plasticity
keeps biological firing rates in a working range \lit \citep{turrigiano2004homeostatic}. These mechanisms are often
introduced together, in one system, and compared against an unmodified baseline. When such a system does better, it
is hard to know which mechanism did the work, or whether a simpler non-biological control would have done the same.

We ask a narrower question: \emph{in a controlled task-incremental continual RL setting, which brain-inspired
mechanisms improve learning or retention when each is isolated and compared with a matched control?} We build one
agent (an SNN or CNN encoder, optional Transformer working memory, PPO) and switch six mechanisms on and off: the
spiking encoder, STDP, homeostasis, direct feedback alignment (DFA) as backprop-free encoder credit assignment (the
heads use exact gradients), sleep-like consolidation, and a shared trunk with per-task heads versus parameter
isolation. Each gets a control chosen to remove the most obvious alternative explanation: a same-size random update
for STDP, a replay baseline with an exactly matched number of replayed samples for sleep, and total-memory accounting
for the shared trunk. A seventh mechanism, Transformer working memory, was to be compared with a memoryless baseline
and a frame stack, but that comparison is inconclusive because benchmark validity was not established.

\paragraph{Headline answer.} \ours No mechanism produced a statistically reliable improvement in final accuracy or
forgetting beyond its matched non-biological control. Given the seed counts, these results rule out large effects
rather than show equivalence, and they rule out large effects only where the intervals are narrow
(Section~\ref{sec:stats}). Two narrower effects stand, and one of them is conditional. Homeostasis speeds up learning
when the spiking encoder uses backprop-free encoder credit assignment (DFA; the heads use exact gradients), without
making it reliable. With a small replay buffer, a shared trunk with per-task heads and sleep reaches similar mean
accuracy to full-width isolation (\smallSnnTwoHundredMean vs \snnIsoMean; \smallSnnTwoHundredN seeds, wide CI) in less
than half the memory; at the default buffer size, isolation is more memory-efficient; and a post-freeze follow-up
found that isolation with quarter-width networks beats shared-trunk methods at equal or lower memory on a 5-task CNN
sequence, so this exception holds only against full-width isolation. The spiking encoder reaches a lower
operation-count regime than a sparsified CNN, at a small accuracy cost; whether that saves energy depends on
neuromorphic hardware we did not use. The Transformer question is inconclusive.

\paragraph{Contributions.}
\begin{enumerate}
  \item A controlled ablation of six brain-inspired mechanisms for task-incremental continual RL, each switched on
        alone and compared with a control chosen to rule out its most obvious alternative explanation, plus a
        Transformer working-memory comparison that remained inconclusive because no benchmark both required memory
        and was learnable within budget (Sections~\ref{sec:mechanisms} and~\ref{sec:results}).
  \item Budget-matched comparisons that change the conclusion of the unmatched ones: sleep versus replay at an
        exactly matched number of replayed samples; a shared trunk with per-task heads versus isolation per unit of
        total memory (parameters plus replay buffer) rather than per parameter; and, after the freeze, isolation at
        a smaller network width (Sections~\ref{sec:res-sleep} and~\ref{sec:res-shared}).
  \item Tests prospectively specified in a version-controlled experiment log, with effect sizes, 95\% confidence
        intervals and multiple-comparison correction, every result labelled as confirmatory, exploratory or
        descriptive (Section~\ref{sec:stats}), and a record of seven earlier positive findings that shrank or
        disappeared under controls, each with its cause (Section~\ref{sec:didnt-survive}).
  \item Code, frozen results, a script that regenerates every number in this paper from the result files, and an
        offline results page with the trained agents (Reproducibility statement).
\end{enumerate}

\paragraph{How to read the claims.} We label statements as literature claims \lit, our interpretations \interp, or
results of our experiments. Each result of ours carries one of three labels. \conf: a test whose comparison,
statistic, seed count and stopping rule were prospectively specified in the version-controlled experiment log before
its runs, and that was run as specified. \expl: a comparison that was not specified in advance (the first two phases
of the project, pilots, context arms, analyses chosen after seeing the data); its $p$-values are uncorrected and
should be read as descriptive. \desc: numbers reported without a test (means, ratios, operation counts, measured CPU
time, memory accounting). \postfz marks the one result obtained after the main results were frozen
(Section~\ref{sec:res-shared}). Summaries that combine several results are marked \ours and refer to the labelled
results. Literature numbers were taken from the cited papers and are marked as such.

% =====================================================================================================
\section{Related work}
\label{sec:related}

\paragraph{Continual RL.} \lit \citet{khetarpal2022towards} review continual RL and its desiderata: retention,
forward transfer, and ideally operation without task identities or boundaries. EWC \citep{kirkpatrick2017overcoming}
penalises moving parameters that were important for earlier tasks (Fisher-weighted); it protects old tasks at the
cost of plasticity for new ones. Learning without Forgetting \citep[LwF;][]{li2017learning} instead regularises the
network's \emph{outputs}: it distils a copy of the network taken before the new task on new-task inputs, so it needs
no stored data. CLEAR \citep{rolnick2019experience} mixes new experience (for plasticity) with
replayed old experience (for stability), trains on both with the off-policy V-trace actor--critic loss, and adds
behavioural cloning on the replayed experience: a KL term towards the stored policy and an L2 term towards the stored
value. It needs no task identities, and it is a strong, simple continual RL method. Our replay baseline and our sleep
phase both reuse CLEAR's cloning losses; Section~\ref{sec:mech-sleep} states exactly how they differ from CLEAR and
from each other.
Continual World \citep{wolczyk2021continual} argues that forgetting alone is an incomplete metric and that forward
transfer must be measured; we report their forward-transfer measure. PackNet \citep{mallya2018packnet} isolates
parameters per task by pruning and freezing, which removes forgetting by construction; our isolation baseline is
the simpler one-network-per-task variant. We measure forgetting as in \citet{chaudhry2018riemannian}.

\paragraph{Sleep and consolidation.} \lit \citet{tadros2022sleep} implement a sleep-like phase in feed-forward
networks (offline, noisy input, local unsupervised Hebbian plasticity) and report that it recovers previously
forgotten tasks in incremental classification. Our sleep phase is \emph{inspired} by this line of work but is
different: it is offline replay plus distillation in an RL setting, optionally with STDP (Section~\ref{sec:mech-sleep}),
and it relies on a replay buffer. \citet{yu2025selfconsolidation} distil accumulated experience of LLM agents into
parameters instead of retrieving raw trajectories, which motivates counting the memory that replay needs
(Section~\ref{sec:res-shared}).

\paragraph{Spiking networks for RL and control.} \lit SNNs are trained end to end with surrogate gradients
\citep{neftci2019surrogate}. \citet{vandenberghe2025adaptive} study the surrogate-gradient slope and show that it
matters much more in RL than in supervised learning; they also describe a warm-up problem for stateful SNNs in RL.
Our encoder resets its membrane at every environment step, so it avoids that problem, and it has no memory across
frames (Section~\ref{sec:mech-snn}). \citet{lee2025brain} and \citet{huo2025research} review biologically plausible
learning rules for SNNs, including STDP and reward-modulated three-factor rules \citep{bi1998synaptic,
fremaux2016neuromodulated}. The Spiking Decision Transformer \citep{pandey2025spiking} combines spiking attention with
three-factor plasticity for offline RL on classic control tasks; Spikingformer \citep{zhou2026spikingformer}
points out that many ``spiking'' architectures hide non-spike floating-point computation, a warning we take into
account when counting operations. SpikingJelly \citep{fang2023spikingjelly} and snnTorch \citep{eshraghian2023training}
are common SNN toolkits; we use our own LIF implementation, checked against snnTorch.

\paragraph{Local and weight-transport-free learning.} \lit Direct feedback alignment \citep{nokland2016direct}
projects the output error to hidden layers through fixed random matrices, avoiding the weight transport that
backpropagation requires. Homeostatic plasticity \citep{turrigiano2004homeostatic} adjusts excitability to keep
firing rates near a set point.

\paragraph{Estimated arithmetic energy proxies.} \lit A common proxy counts synaptic operations and charges an
accumulate (about \pjAc~pJ) per spike-driven operation and a multiply--accumulate (about \pjMac~pJ) per dense
operation, using 45\,nm estimates \citep{horowitz2014computing}. Such numbers are estimates of arithmetic energy, not
measured power, and they assume event-driven hardware.

% =====================================================================================================
\section{Setup}
\label{sec:setup}

\subsection{Environments and task sequences}
All environments are MiniGrid \citep{chevalierboisvert2023minigrid} gridworlds with the standard $7{\times}7$
egocentric symbolic view, which we encode as 20 one-hot channels (object type, colour and state).

\paragraph{fetch3 and fetch5 (continual).} A room ($6{\times}6$ interior) contains three (fetch3) or five (fetch5)
distinct objects at random positions. Task $k$ rewards picking up object $k$ (reward $1-0.9\,t/T_{\max}$, the MiniGrid
convention); picking up any other object ends the episode with reward 0. Observations are statistically identical
across tasks, so the same state demands different actions in different tasks: the setting guarantees interference
under naive fine-tuning, while navigation and pick-up are shared skills that allow forward transfer. The agent is
given the task identity (a task embedding and one actor--critic head per task). This is the standard
\emph{task-incremental} setting; without the task identity the tasks are not merely harder but unsolvable, since the
goal is not visible (Section~\ref{sec:limitations}). Tasks are trained in order for a fixed number of frames each
(\framesPerTaskFlat per task for the CNN and SNN agents, \framesPerTaskHybrid for the SNN+Transformer agent).

\paragraph{DoorKey-6x6 (single task).} Used for the local-plasticity questions: the agent must pick up a key, open
a door and reach the goal; \framesDoorKey frames. It is hard enough that the encoder's learning matters (a frozen
random encoder learns more slowly), and outcomes are bimodal (a run either solves it or stays near zero).

\paragraph{MiniGrid-Memory (S7, S11, S13).} The agent sees a cue object at one end of a corridor and must choose the
matching object at the other end. We use these maps for the Transformer question (Section~\ref{sec:res-memory}).

\subsection{Agent and training}
\label{sec:agent}
The encoder is either a CNN (three $2{\times}2$ convolutions with \convChannels channels and a linear layer to
\featDim features) or an SNN with the same topology, leaky integrate-and-fire (LIF) neurons, \snnT timesteps per
environment step, a learnable leak (initial $\beta=\lifBeta$) and a surrogate gradient (slope \surrogateSlope)
\citep{neftci2019surrogate}. Actor and critic heads are small MLPs. Training is PPO \citep{schulman2017proximal}
in a single-file CleanRL-style implementation \citep{huang2022cleanrl}: \numEnvs parallel environments,
\numSteps steps per rollout, \updateEpochs epochs, minibatches of \minibatch, clip \clipCoef, $\gamma=\gammaVal$,
GAE $\lambda=\gaeLambda$, early stopping at an approximate KL of \targetKl, Adam with learning rate \lrFlat
(\lrHybrid for the SNN+Transformer agent). Everything runs on CPU (one thread per run).

\paragraph{Shared trunk with per-task heads.} Every non-isolation agent is a shared trunk with per-task heads. The
trunk (the encoder, plus the Transformer in the SNN+Transformer agent) is shared by all tasks; each task adds one
actor--critic head and one task-embedding vector. For the CNN agent the trunk has \paramsCnnTrunk parameters and each
task adds \paramsPerTaskShared (a \paramsHead-parameter head and a \paramsTaskEmb-dimensional embedding), so the
3-task agent has \paramsCnnShared parameters; the SNN and SNN+Transformer agents differ only in the trunk (the shared
SNN+Transformer agent has \paramsHybridShared). The parameters that scale with the number of tasks are therefore the
heads and embeddings, and, for replay and sleep, the replay buffer (\bufferPerTask stored states per task by default).
Isolation instead adds a full network of \paramsCnnPerTask parameters (CNN) per task.

\subsection{Baselines: three strategies against forgetting}
\label{sec:baselines}
The non-biological comparators follow the three standard strategies \lit \citep{khetarpal2022towards}.
\emph{Regularisation}: EWC \citep{kirkpatrick2017overcoming} penalises changes to parameters that were important for
earlier tasks; our online variant keeps one accumulated Fisher estimate and one anchor copy of the parameters, so its
storage does not grow with the number of tasks ($\lambda=10,100,1000$). \emph{Replay}: CLEAR-style replay
\citep{rolnick2019experience} interleaves cloning losses on stored states with learning the new task
(Section~\ref{sec:mech-sleep}); it stores a buffer that grows with the number of tasks. \emph{Capacity isolation}: one
fresh network per task, which cannot forget by construction and whose parameters grow by a full network per task.
Naive fine-tuning of the shared trunk is the lower bound. Sleep-like consolidation is compared with replay because it
uses the same buffer and losses. The post-freeze follow-up (Section~\ref{sec:res-shared}) adds Learning without
Forgetting \citep[LwF;][]{li2017learning}, a functional regulariser: at the start of each task it keeps a frozen copy of
the network (stored in int8, \postSnapshotMb~MB for the 5-task CNN agent) and, on every PPO minibatch, distils its
policy and value on \postLwfBatch current-task states, each paired with a random earlier task's head. It stores no
data between tasks.

\subsection{Metrics}
$R_{i,j}$ is the mean evaluation return on task $j$ after training on task $i$ (\evalEpisodesCl episodes, actions
sampled from the policy). For $T$ tasks: ACC $=\frac1T\sum_j R_{T,j}$ (final average performance); BWT
$=\frac{1}{T-1}\sum_{j<T} (R_{T,j}-R_{j,j})$; FORGET $=\frac{1}{T-1}\sum_{j<T}\big(\max_{i<T}R_{i,j}-R_{T,j}\big)$
\citep{chaudhry2018riemannian}; forward transfer FWT is the normalised gain in area under the training curve on
tasks $2..T$ relative to a from-scratch network on the same task and seed \citep{wolczyk2021continual}. For single-task
runs we report the solve rate (final evaluation return above \solveThreshold over \evalEpisodesSingle episodes) and
the AUC (mean training return over the run, a sample-efficiency measure).

\paragraph{Memory accounting.} For the shared-trunk question we count total memory as fp32 parameters plus the
replay buffer at its measured storage cost: \bytesPerState bytes per stored state for the CNN and SNN agents
(observation, mask, task id, teacher logits and value) and \bytesPerStateHybrid bytes for the SNN+Transformer agent
(4-frame windows), times \bufferPerTask states per task by default.

\subsection{Statistics}
\label{sec:stats}
Each arm has \seedRange seeds. For solve rates we use Fisher's exact test and the Newcombe (Wilson score) interval for the
difference of proportions \citep{newcombe1998interval}; for continuous metrics, Welch's $t$-test
\citep{welch1947generalization} and the Welch 95\% confidence interval of the difference of means. Where a
prospectively specified family contains several contrasts we apply Holm's correction \citep{holm1979simple}. With so
few seeds, CIs are wide; we treat them as the primary result. A non-significant difference is read as ``no large
effect'' only where its interval excludes large effects, and never as equivalence.

\paragraph{Seed sets.} Within every comparison, all arms use the same seed set: seeds 1 to $n$ for the frozen runs and
\postSeeds for the post-freeze follow-up, so each seed appears in every arm of the comparison. The one exception is the
unmatched SNN+Transformer sleep arm (seeds 1--\nHybSleep), which Table~\ref{tab:continual} and
Figure~\ref{fig:forgetting} show next to replay on seeds 1--\nHybReplay; the budget-matched test
(Section~\ref{sec:res-sleep}) uses seeds 1--\matchN in both arms, and the seed-count history in
Table~\ref{tab:survive} restricts both arms to the same seeds. The tests treat the arms as independent samples.

\paragraph{Prospective specification.} The work was done in four phases. The first two were exploratory. For the last
two, the comparisons, test statistics, seed counts and stopping rules were recorded in the project's
version-controlled experiment log before the runs, one planned test per question; the post-freeze follow-up was
specified in the same way in its own log (\texttt{explore/EXPLORE.md}) before its runs. This is a project record, not
a public registry. We call these tests \emph{prospectively specified in a version-controlled experiment log}
(prospectively specified, for short) and label their results \conf. Results not specified in advance are labelled
\expl and enter no confirmatory test.

% =====================================================================================================
\section{Mechanisms and their controls}
\label{sec:mechanisms}

Table~\ref{tab:mechanisms} lists each mechanism, how it is implemented and what it is compared against.

\begin{table}[t]
\centering
\small
\caption{Mechanisms, implementations and controls. Each row is switched on and off independently. The first six rows
were tested; the Transformer comparison is inconclusive because benchmark validity was not established.}
\label{tab:mechanisms}
\renewcommand{\arraystretch}{1.25}
\begin{tabular}{>{\raggedright\arraybackslash}p{2.6cm}>{\raggedright\arraybackslash}p{6.2cm}>{\raggedright\arraybackslash}p{4.4cm}}
\toprule
Mechanism & Implementation & Control / comparator \\
\midrule
Spiking encoder & LIF CNN, \snnT timesteps per frame, surrogate gradients; membrane reset every environment step & Same-topology CNN; op count at matched accuracy \\
STDP & Reward-modulated three-factor STDP on encoder synapses, added to the backprop update: $\Delta w=\alpha\,\Delta w_{\text{STDP}}+\Delta w_{\text{grad}}$, stabilised by mean-centring and homeostasis & Same-RMS random update + homeostasis; homeostasis alone \\
Homeostasis & Per-neuron bias pushes firing rate towards \homeoTarget & Same agent without it \\
DFA & Backprop-free encoder credit assignment: encoder layers learn from the output error through fixed random feedback (heads use exact gradients) & DFA without homeostasis; backprop \\
Sleep & Offline phases: old-task distillation on stored states + current-task self-distillation & CLEAR-style replay (cloning losses, interleaved), replayed samples matched exactly; EWC; naive \\
Shared trunk & One trunk shared by all tasks, per-task heads and task embeddings & One network per task (isolation), at equal total memory; after the freeze, isolation at quarter width \\
\midrule
Transformer memory (inconclusive) & Transformer over the last $K$ encoded frames (past frames only) & Memoryless agent; frame stack of $K$ frames \\
\bottomrule
\end{tabular}
\end{table}

\subsection{Spiking encoder}
\label{sec:mech-snn}
The SNN integrates each frame over \snnT internal timesteps and its membrane state is reset at every environment
step, so it has no memory across frames; we verified this empirically (identical output for identical frames
regardless of history). We count event-driven operations per frame: synaptic operations for the SNN (non-zero
presynaptic spikes times fan-out, summed over timesteps) and non-zero-activation multiply--accumulates for a CNN
whose activations are sparsified with an L1 penalty. The estimated arithmetic energy proxy multiplies these counts by
the literature constants above. To compare accuracy against operations without RL noise, the op-count frontier is
measured by behaviour cloning of a trained DoorKey-6x6 teacher: both encoders plus an identical linear readout are
trained with an activity penalty swept over a range, and we report held-out action accuracy and the closed-loop
return of the cloned policy (\rqThreeSeeds seeds per point).

\subsection{STDP and the random-update control}
\label{sec:mech-stdp}
After every PPO optimiser step, spikes recorded in the minibatch's forward pass produce a local STDP update for each
encoder layer (eligibility traces, pre/post pairing), multiplied by a third factor, the normalised PPO advantage of
the sample (reward-modulated STDP) \citep{fremaux2016neuromodulated, lee2025brain}. The update is normalised to unit
RMS and scaled by $\alpha$ times the learning rate, so $\alpha=1$ moves each weight about as much as one Adam step.
The \emph{stabilised} version mean-centres pre- and post-synaptic spikes (a covariance rule, no drift under
independent firing) and adds homeostasis. The \emph{random} control applies an update with the same per-layer RMS
but a random direction, also with homeostasis, so it separates ``STDP helps'' from ``any perturbation of this size
helps''. We use $\alpha=\stdpAlpha$ for the prospectively specified comparison.

\subsection{Homeostasis}
Each spiking neuron's bias is nudged towards a target firing rate: $b\leftarrow b+\eta\,(\text{\homeoTarget}-r)$ with
$\eta=\homeoLr$ and $r$ the neuron's rate in the minibatch.

\subsection{Direct feedback alignment}
In the DFA agents, credit assignment in the encoder is backprop-free: the encoder's layers receive the output error
through fixed random feedback matrices instead of backpropagated gradients \citep{nokland2016direct}, while the actor
and critic heads use exact gradients. This gives the encoder a global learning signal without weight transport.

\subsection{Sleep-like consolidation and the replay baseline}
\label{sec:mech-sleep}
Every \sleepPeriod PPO updates, and once at the end of each task, the agent goes offline for \sleepSteps gradient
steps on two losses: (i) distillation towards the stored policy and value on states saved at the end of each
earlier task (KL divergence plus value MSE; the same buffer the replay baseline uses), and (ii) self-distillation of
the current task towards a frozen snapshot of the network taken at sleep onset, on a reservoir of recent wake states.
Optionally, local STDP runs on the spikes produced during sleep. The \emph{replay} baseline adds the same distillation
loss on one replayed minibatch to every PPO minibatch while learning later tasks, as in CLEAR.

\paragraph{Relation to CLEAR.} Both methods reuse CLEAR's behavioural-cloning losses \citep{rolnick2019experience},
and they differ from CLEAR, and from each other, in specific ways. (i) Our \emph{replay} baseline keeps CLEAR's
cloning terms (policy KL and value regression on replayed states) but not its off-policy RL loss on replayed
experience: PPO is on-policy, and replayed transitions cannot enter its objective without an importance-weighting
scheme such as V-trace, which we did not add. The cloning targets are the policy and value of the network at the end
of the task a state comes from, stored with the state, rather than those of the policy that generated the experience.
As in CLEAR, replay is interleaved with every update on the new task. (ii) Our \emph{sleep} phase uses the same buffer
and the same cloning losses, but applies them offline, in scheduled phases without environment interaction, and adds
self-distillation on recent current-task states. With the number of replayed states matched, the sleep--replay
comparison therefore isolates \emph{when} replay happens, plus the self-distillation term; it does not compare sleep
with the full CLEAR method. (iii) Unlike CLEAR, both use the task identity (per-task heads), because our setting is
task-incremental.

\paragraph{Replay budget.} The two methods consume different numbers of replayed states. Replay's count depends on
how many PPO minibatches run, which varies by seed because of KL early stopping (\matchSamplesMin--\matchSamplesMax
per SNN+Transformer run); sleep's is fixed by its schedule (\unmatchedSleepSamples). For the prospectively specified
comparison, each sleep run was given a total replay budget equal to the measured count of the same-seed replay run,
spread evenly over its sleep phases (\emph{sleep, replay-matched}); the match was exact for every seed.

\subsection{Shared trunk with per-task heads versus isolation}
The shared agent keeps one trunk and adds a head and a task embedding per task (Section~\ref{sec:agent}). Isolation
trains a fresh network per task: zero forgetting by construction and a full network's parameters per task. We compare
them per parameter, as is common, and per unit of total memory, which includes the replay buffer that the
shared-trunk methods need. In the frozen study, each isolated network has the same width as the shared trunk; the
post-freeze follow-up also tests a narrower isolated network.

\subsection{Transformer working memory}
The SNN+Transformer agent encodes each of the last $K$ frames and runs a one-layer Transformer encoder
\citep{vaswani2017attention} over them, reading out the newest position; the window holds only past frames, and padding
before the episode start is masked. On the fetch tasks $K=\hybridWindow$.
The non-attention baselines are a frame stack (the last $K$ frames as input channels) and, implemented but not run, a
GRU over the same encoded frames.

% =====================================================================================================
\section{Results}
\label{sec:results}

\begin{figure}[t]
\centering
\fig{fig_frontier}
\caption{Accuracy against operation count for the spiking and the conventional encoder (descriptive). Each encoder,
with an identical linear readout, is trained by behaviour cloning of a trained DoorKey-6x6 PPO teacher with an activity
penalty; each point is one penalty value (mean of \rqThreeSeeds seeds, the same seeds at every point). The $x$-axis
counts event-driven operations per frame (synaptic operations for the SNN, non-zero multiply--accumulates for the
L1-sparsified CNN; log scale). Left: held-out action accuracy (dashed: majority-action rate). Right: closed-loop return
of the cloned policy. The SNN keeps working below the lowest operation count at which the CNN still works.}
\label{fig:frontier}
\end{figure}

\subsection{Spiking encoder: a lower operation-count regime, not more accuracy}
\label{sec:res-snn}

\desc Where both encoders work, they are equally accurate at matched operation counts (Figure~\ref{fig:frontier}).
The unpenalised CNN uses \rqThreeCnnDenseOps event-driven operations per frame (accuracy \rqThreeCnnDenseAcc).
The L1-sparsified CNN still works at \rqThreeCnnFloorOps operations (accuracy \rqThreeCnnFloorAcc), but a stronger
penalty silences the whole network at once (\rqThreeCnnCollapseOps operations, accuracy \rqThreeCnnCollapseAcc,
the majority-action rate). The SNN keeps working below that floor: with \snnT timesteps reduced to 2 it reaches
\rqThreeSnnTwoLowOps operations at accuracy \rqThreeSnnTwoLowAcc and a closed-loop return of
\rqThreeSnnTwoLowRet, about \rqThreeOpsRatio{}$\times$ fewer operations than the CNN floor; $T=4$ behaves similarly
(\rqThreeSnnFourLowOps, \rqThreeSnnFourLowAcc). Under the estimated arithmetic energy proxy (literature constants for
45\,nm arithmetic), this point costs \rqThreeSnnTwoLowEnergy~nJ per frame against \rqThreeCnnFloorEnergy~nJ for the
sparsest working CNN and \rqThreeCnnDenseEnergy~nJ for the dense CNN.

\desc Measured cost in our implementation: on the CPU we actually used, one PPO training step takes \msSnn~ms for the
SNN against \msCnn~ms for the CNN (\snnCpuSlowdown{}$\times$), because the SNN simulates \snnT timesteps. The unit is
CPU process time on one thread for one forward pass, backward pass and Adam update on a minibatch of 256 DoorKey-6x6
observations (single-task networks), averaged over 5 repetitions after one warm-up step (\texttt{scripts/benchmark\_compute.py}).
\interp The SNN's advantage is access to a lower-operation regime, not better accuracy, and the energy numbers assume
event-driven hardware. Other CNN sparsifiers (e.g.\ k-winner-take-all) might move the CNN floor; we did not test them.

\paragraph{Robustness per unit of compute.} \desc Dividing final continual accuracy (fetch3) by training compute (the
measured per-step CPU cost above times frames), the plain CNN scores \rqFourCnn per unit, \rqFourRatio{}$\times$ the
best spiking variant (SNN with sleep, \rqFourSnn); SNN with homeostasis scores \rqFourSnnHomeo, DFA with homeostasis
\rqFourDfaHomeo, stabilised STDP \rqFourStdp and the SNN+Transformer agent \rqFourHybrid. Encoders trained by local
STDP only never solved DoorKey (\rqFourLocalOnly runs).

\begin{figure}[t]
\centering
\fig{fig_rq1}
\caption{Local-plasticity arms on DoorKey-6x6 (confirmatory: prospectively specified; \framesDoorKey frames,
\rqOneBpN seeds per arm, seeds 1--\rqOneBpN in every arm). Each dot is one seed's AUC (mean training return over the
run); filled dots solved the task (final evaluation return above \solveThreshold), open dots did not. Black bars: mean
with 95\% CI. Numbers above the dots: seeds solved. Left group: backpropagation as the global learning signal, alone or
with homeostasis, stabilised reward-modulated STDP, or a random update of the same size. Right group: backprop-free
encoder credit assignment by direct feedback alignment (DFA; heads use exact gradients), with and without homeostasis.}
\label{fig:rq1}
\end{figure}

\begin{table}[t]
\centering
\small
\caption{Prospectively specified local-plasticity contrasts on DoorKey-6x6 (confirmatory; \rqOneBpN seeds per arm,
seeds 1--\rqOneBpN in every arm). Differences are the first arm minus the second, with 95\% CIs (Newcombe for solve
rates, Welch for AUC). Holm-corrected $p$ across the four contrasts, separately per statistic.}
\label{tab:rq1}
\footnotesize
\setlength{\tabcolsep}{4pt}
\begin{tabular}{lcccc}
\toprule
Contrast & $\Delta$ solve rate [95\% CI] & Fisher $p$ (Holm) & $\Delta$ AUC [95\% CI] & Welch $p$ (Holm) \\
\midrule
STDP vs homeostasis alone & \rqOneContrASolveCi & \rqOneContrAFisherP~(\rqOneContrAFisherHolm) & \rqOneContrAAucCi & \rqOneContrAAucP~(\rqOneContrAAucHolm) \\
STDP vs matched random & \rqOneContrBSolveCi & \rqOneContrBFisherP~(\rqOneContrBFisherHolm) & \rqOneContrBAucCi & \rqOneContrBAucP~(\rqOneContrBAucHolm) \\
homeostasis vs backprop & \rqOneContrCSolveCi & \rqOneContrCFisherP~(\rqOneContrCFisherHolm) & \rqOneContrCAucCi & \rqOneContrCAucP~(\rqOneContrCAucHolm) \\
random vs homeostasis & \rqOneContrDSolveCi & \rqOneContrDFisherP~(\rqOneContrDFisherHolm) & \rqOneContrDAucCi & \rqOneContrDAucP~(\rqOneContrDAucHolm) \\
\bottomrule
\end{tabular}
\end{table}

\subsection{STDP: the tested rule was not distinguishable from a same-size random update}
\label{sec:res-stdp}

\conf With backpropagation as the global signal, the four arms solve DoorKey-6x6 at similar rates: backprop
\rqOneBpSolved, backprop with homeostasis \rqOneHomeoSolved, with stabilised STDP \rqOneStdpSolved, with the matched
random update \rqOneRandSolved (Figure~\ref{fig:rq1}). All four prospectively specified contrasts are null after Holm
correction (Table~\ref{tab:rq1}). The tested rule (stabilised reward-modulated STDP at $\alpha=\stdpAlpha$, added to
the backprop update) was not distinguishable from a random update of the same size (AUC difference
\rqOneContrBAucCi, $p=\rqOneContrBAucP$), and its point estimate was slightly below homeostasis alone. These intervals
are wide: with \rqOneBpN seeds per arm they do not exclude a large effect in either direction, so this is ``no
detectable effect of this rule in this setup'', not evidence that STDP cannot help.

\expl Unstabilised STDP hurt. In an earlier DoorKey-5x5 experiment, reward-modulated STDP with $\alpha=1$ solved
\vanillaStdpSolved runs against \vanillaBpSolved for backprop (AUC difference \vanillaStdpAucCi, $p=\vanillaStdpP$),
with a mean spike rate of \vanillaStdpRate over the last updates against \bpRate for backprop. A larger stabilised
$\alpha$ (\stdpAlphaLarge) was also worse (\rqOneStdpLargeSolved solved, context arm).
\interp The frame is static across the SNN's timesteps, so the same pre--post pairings repeat and a mostly
potentiating rule feeds back on itself (weights up, rates up, more correlated spikes). Stabilisation removes the
runaway but leaves nothing measurable beyond a random perturbation.

\expl Inside sleep, the same pattern holds on the SNN trunk (fetch3): vanilla STDP during sleep collapsed performance
(ACC \snnSleepVanillaStdpAcc, \snnSleepVanillaStdpN seed), while stabilised STDP (\snnSleepStabStdpAcc) was
indistinguishable from plain sleep (\accSnnSleep).

\subsection{Homeostasis and DFA: faster learning with backprop-free encoder credit assignment, not reliable learning}
\label{sec:res-homeo}

\conf With backprop, homeostasis raised the AUC by \rqOneContrCAucCi (raw $p=\rqOneContrCAucP$), which does not
survive Holm correction ($p=\rqOneContrCAucHolm$); its solve rate was \rqOneHomeoSolved vs \rqOneBpSolved
($p=\rqOneContrCFisherP$).

\conf With backprop-free encoder credit assignment by DFA (heads use exact gradients; prospectively specified,
\rqOneDfaN seeds per arm, seeds 1--\rqOneDfaN in both arms), homeostasis raised the solve rate from \rqOneDfaSolved to
\rqOneDfaHomeoSolved (difference \dfaSolveCi; Fisher $p=\dfaFisherP$, Holm \dfaFisherHolm, not significant) and the
AUC from \rqOneDfaAuc to \rqOneDfaHomeoAuc (difference \dfaAucCi, \dfaAucRatio{}$\times$; Welch $p=\dfaAucP$, Holm
\dfaAucHolm, significant). The remaining homeostasis seeds still failed outright (Figure~\ref{fig:rq1}).
\interp Homeostasis is a real but partial fix for DFA: it speeds up learning without a weight-transport requirement in
the encoder, but it does not make DFA reliable on this task, and backprop with homeostasis (\rqOneHomeoSolved) remains
better. \expl In continual learning, a DFA+homeostasis encoder with sleep reached only \dfaContinualAcc ACC on fetch3.

\begin{figure}[t]
\centering
\fig{fig_forgetting}
\caption{Forgetting on the continual task sequences, by method (exploratory, except the replay-matched sleep arm, which
belongs to the confirmatory test of Figure~\ref{fig:budget}). Columns: agent and sequence. Top row: return on the
first task after training on each task in turn (a flat line means task 1 is not forgotten). Bottom row: mean return over
the tasks seen so far. Mean over seeds with 95\% CIs (\nCnnSleep seeds per arm, seeds 1--\nCnnSleep, except the
SNN+Transformer agent: replay seeds 1--\nHybReplay, replay-matched sleep 1--\nHybSleepMatched, sleep 1--\nHybSleep); the
CIs are clipped to $[0,1]$, the range of returns, and with \nCnnSleep seeds the SNN intervals span most of it. Methods
are offset horizontally so that overlapping lines stay visible. Colours mean the same method in every figure.}
\label{fig:forgetting}
\end{figure}

\begin{table}[t]
\centering
\small
\caption{Final ACC (mean $\pm$ SD over seeds) and FORGET on the continual sequences (exploratory). ``Sleep, matched''
consumes exactly as many replayed states as the same-seed replay run (SNN+Transformer only; confirmatory test in
Section~\ref{sec:res-sleep}). Seeds 1--\nCnnSleep in every arm, except the SNN+Transformer column (replay and sleep,
matched: seeds 1--\nHybReplay; sleep: 1--\nHybSleep; naive and isolation: 1--\nHybNaive).}
\label{tab:continual}
\footnotesize
\setlength{\tabcolsep}{3.5pt}
\begin{tabular}{lcccc}
\toprule
& CNN, 3 tasks & SNN, 3 tasks & SNN+Transformer, 3 tasks & CNN, 5 tasks \\
\midrule
\multicolumn{5}{l}{\emph{Final ACC}} \\
naive & \accCnnNaive & \accSnnNaive & \accHybNaive & \accCnnFiveNaive \\
EWC ($\lambda{=}100$) & \accCnnEwc & \accSnnEwc & -- & -- \\
replay & \accCnnReplay & \accSnnReplay & \accHybReplay & \accCnnFiveReplay \\
sleep & \accCnnSleep & \accSnnSleep & \accHybSleep & \accCnnFiveSleep \\
sleep, matched & -- & -- & \accHybSleepMatched & -- \\
isolation & \accCnnIso & \accSnnIso & \accHybIso & \accCnnFiveIso \\
\midrule
\multicolumn{5}{l}{\emph{FORGET (mean)}} \\
naive & \forgetCnnNaive & \forgetSnnNaive & \forgetHybNaive & \forgetCnnFiveNaive \\
replay & \forgetCnnReplay & \forgetSnnReplay & \forgetHybReplay & \forgetCnnFiveReplay \\
sleep & \forgetCnnSleep & \forgetSnnSleep & \forgetHybSleep & \forgetCnnFiveSleep \\
sleep, matched & -- & -- & \forgetHybSleepMatched & -- \\
\midrule
\multicolumn{5}{l}{\emph{sleep $-$ replay, final ACC [95\% Welch CI] (unmatched budget)}} \\
& \sleepVsReplayCnnCi & \sleepVsReplaySnnCi & \sleepVsReplayHybCi & \sleepVsReplayCnnFiveCi \\
\bottomrule
\end{tabular}
\end{table}

\subsection{Sleep-like consolidation: no detectable difference from replay at a matched replay budget}
\label{sec:res-sleep}

\expl Sleep and replay both remove most of naive fine-tuning's forgetting on every agent (Figure~\ref{fig:forgetting},
Table~\ref{tab:continual}); EWC protects old tasks at a plasticity cost that grows with $\lambda$ (return on the last
task right after training it: \ewcLastTaskTen, \ewcLastTaskHundred and \ewcLastTaskThousand for $\lambda=10,100,1000$,
against \naiveLastTaskCnn for naive; CNN). At unmatched replay budgets, sleep never differs significantly from replay
(last row of Table~\ref{tab:continual}).

\begin{figure}[t]
\centering
\fig{fig_budget}
\caption{Sleep versus replay at an exactly matched replay budget (SNN+Transformer agent, fetch3; confirmatory for final
ACC, prospectively specified; \matchN seeds per arm, seeds 1--\matchN in both arms). \emph{Replay}: CLEAR-style replay.
\emph{Sleep (matched)}: each run replays exactly as many stored states as the replay run with the same seed.
\emph{Sleep (+\unmatchedExtraPct\%)}: the original sleep runs (seeds 1--\nHybSleep), which replayed \unmatchedExtraPct\%
more states. Grey lines connect runs with the same seed; black bars are means with 95\% CIs. Left: final accuracy
(ACC). Right: forgetting (FORGET; lower is better; descriptive).}
\label{fig:budget}
\end{figure}

\conf \textbf{Replay budget matched.} On the SNN+Transformer agent, where the original sleep runs had consumed
\unmatchedExtraPct\% more replayed states than replay, sleep with an exactly matched budget was not distinguishable
from replay on final accuracy (\matchAccSleep vs \matchAccReplay; difference \matchAccCi, $p=\matchAccP$, \matchN seeds
per arm, the same seeds in both arms). Unlike the STDP contrasts, this interval is narrow. \desc On forgetting, specified as descriptive in the log, the arms were also close (\matchForgetSleep vs
\matchForgetReplay; difference \matchForgetCi, $p=\matchForgetP$) (Figure~\ref{fig:budget}). On the same seeds, the
unmatched sleep runs forgot \unmatchedForgetSameSeeds on average and the matched ones \matchedForgetSameSeeds; matched
minus unmatched ACC was \matchedMinusUnmatchedAccCi (paired, $p=\matchedMinusUnmatchedAccP$).
\interp Sleep's small earlier edge on this agent came from its extra replay, not from consolidating offline. This
matches the literature point that replay is a very strong baseline \citep{rolnick2019experience}: in our implementation
sleep is a way of scheduling replay, and at an equal replay budget we found no advantage for it. This is one sleep
schedule, on one agent and sequence, at the full replay budget; it does not show that offline consolidation cannot
help elsewhere.

\begin{figure}[t]
\centering
\fig{fig_memory}
\caption{Final accuracy (ACC) against total memory (descriptive; the SNN comparisons with isolation are confirmatory,
Section~\ref{sec:res-shared}): fp32 parameters plus the replay buffer at its measured storage cost, log scale. Small
dots: seeds; large markers: means. Numbers above the $x$-axis: replay states stored per task by sleep and replay;
isolation stores none but trains one full-width network per task. At the default \bufferPerTask states per task,
isolation needs the least memory; with 200 states per task, sleep reaches a similar mean accuracy in less memory
(\smallSnnTwoHundredN seeds; wide CI). The 5-task SNN runs use the fair budget of \framesPerTaskFiveFair frames per task
(Section~\ref{sec:didnt-survive}). The post-freeze quarter-width isolation result (Section~\ref{sec:res-shared}) is not
plotted.}
\label{fig:memory}
\end{figure}

\subsection{Shared trunk with per-task heads versus isolation: memory-efficient only with small buffers and full-width isolation}
\label{sec:res-shared}

\desc Per parameter, the shared trunk with sleep is clearly more efficient than isolation: ACC per 100k parameters
\perParamCnnThreeSleep vs \perParamCnnThreeIso (CNN), \perParamSnnThreeSleep vs \perParamSnnThreeIso (SNN),
\perParamHybThreeSleep vs \perParamHybThreeIso (SNN+Transformer), \perParamCnnFiveSleep vs \perParamCnnFiveIso (CNN,
5 tasks). But at the default \bufferPerTask stored states per task the buffer outweighs the extra networks: isolation
uses less total memory and gets more accuracy per MB everywhere (\memCnnThreeIsoPerMb vs \memCnnThreeSleepPerMb,
\memSnnThreeIsoPerMb vs \memSnnThreeSleepPerMb, \memHybThreeIsoPerMb vs \memHybThreeSleepPerMb, \memCnnFiveIsoPerMb vs
\memCnnFiveSleepPerMb per MB; Figure~\ref{fig:memory}).

\desc With small buffers the shared trunk can be the more memory-efficient one, against full-width isolation. On the
CNN, sleep with 200 states per task reaches \smallCnnTwoHundredAcc in \smallCnnTwoHundredMb~MB
(\smallCnnTwoHundredPerMb per MB, against \memCnnThreeIsoPerMb for isolation). \conf On the SNN, sleep with 200 states
per task reaches a similar mean accuracy to isolation (\smallSnnTwoHundredAcc vs \snnIsoAcc; difference
\snnSleepTwoHundredVsIsoCi, $p=\snnSleepTwoHundredVsIsoP$, \smallSnnTwoHundredN seeds per arm, the same seeds in both
arms) in \smallSnnTwoHundredMb~MB against \memSnnThreeIsoMb~MB, \smallSnnTwoHundredMemPct\% of the memory. With
\smallSnnTwoHundredN seeds the interval is wide: the data rule out neither a large deficit nor a large advantage, so
this is not evidence of equal accuracy. At \emph{equal} total memory (SNN, 2000 states per task,
\smallSnnTwoThousandMb vs \memSnnThreeIsoMb~MB) sleep does not beat isolation (difference
\snnSleepTwoThousandVsIsoCi, $p=\snnSleepTwoThousandVsIsoP$).

\expl Forward transfer depends on the encoder: positive where learning from scratch is slow (SNN \fwtSnn,
$p=\fwtSnnP$; SNN+Transformer \fwtHyb, $p=\fwtHybP$), near zero on the 3-task CNN (\fwtCnn) and negative on the 5-task
CNN (\fwtCnnFive). \conf On the 5-task SNN at the fair budget it was \fiveSnnFairFwt ($p=\fiveSnnFairFwtP$).

\paragraph{Post-freeze follow-up: isolation's network width.} \postfz\conf The comparisons above use isolated
networks of the same width as the shared trunk (\paramsCnnPerTask parameters each). After the main results were
frozen, we tested whether narrower isolated networks change the memory comparison, in a follow-up whose hypotheses,
test, seed count and stopping rule were prospectively specified in its experiment log (\texttt{explore/EXPLORE.md},
comparison C4) before its runs. On the 5-task CNN sequence, isolation with quarter-width networks (convolution channels
\postNarrowChannels, \postNarrowFeat features, \postNarrowHidden hidden units; \postNarrowPerTask parameters per task,
\postNarrowParams in total, \postNarrowMb~MB) reached a final ACC of \postNarrowAcc. A shared trunk with per-task heads
trained with LwF and an int8 teacher stores only its parameters (\postLwfParams, \postLwfMb~MB, plus a transient
\postSnapshotMb~MB teacher during training); it reached \postLwfAcc (difference \postHEightCi, Holm
$p=\postHEightHolm$). The same trunk with replay (\postReplayBuf stored states per task, \postReplayMb~MB) reached
\postReplayAcc (difference \postHNineCi, Holm $p=\postHNineHolm$). Full-width isolation (\postIsoMb~MB) reached
\postIsoAcc; quarter width minus full width was \postHTenCi (Holm $p=\postHTenHolm$), inside the prospectively specified
$\pm0.05$ margin for ``no large difference''. All four arms used seeds \postSeeds (\postN per arm) and the frozen CNN
protocol (\framesPerTaskFlat frames per task, learning rate \lrFlat, \evalEpisodesCl evaluation episodes per task); LwF
and replay processed the same number of distilled or replayed states (\postSamples per run).
\interp On this sequence the frozen isolation baseline's per-task networks were much larger than the tasks need, so
``a shared trunk is more memory-efficient than isolation with small buffers'' holds only against full-width
isolation. Two limits apply. The shared trunk was not shrunk in the same way, so this is not a width-matched
comparison of the two strategies, and the follow-up used LwF and replay rather than sleep. The SNN case is untested: the
only narrow SNN isolation runs were a two-seed pilot at \framesPerTaskFlat frames per task, a budget at which even
full-width fresh SNNs often had not learned their tasks (Section~\ref{sec:didnt-survive}).

\interp The shared trunk's advantage in the frozen study is memory efficiency when replay is cheap and isolation is
full width, not higher accuracy at equal memory. ``Shared weights are more efficient than isolation'' is only true with
the buffer size and the isolation width stated.

\begin{figure}[t]
\centering
\fig{fig_probe}
\caption{Memory-benchmark probes (sizing runs whose decision rule was prospectively specified; CNN encoder, one seed
each, learning rate \probeLr). Training return (moving mean over 25 updates) on MiniGrid-MemoryS11 and S13 with a single
frame (grey) or a stack of the last \memWindow frames (violet); light violet: an exploratory \probeLongFrames-frame run
of the stacked CNN on S11. A return of 0.5 is chance at the junction. No run learns to use the cue, so neither map tests
working memory.}
\label{fig:probe}
\end{figure}

\subsection{Transformer working memory: inconclusive, because benchmark validity was not established}
\label{sec:res-memory}

\expl On the fetch tasks, which are fully determined by the current frame, the Transformer has nothing to remember;
the SNN+Transformer agent needed \framesPerTaskHybrid frames per task where the CNN and SNN agents used
\framesPerTaskFlat. To test memory we used MiniGrid-Memory. On S7 (pilot, one seed, \sSevenFrames frames) a
memoryless CNN already succeeded in \sSevenCnnOne of episodes (a 4-frame stack: \sSevenCnnFour), so S7 does not
isolate memory: in that small map the cue and both choices are visible from one position. The pilot's SNN runs
(SNN with an 8-frame stack \sSevenSnnEight, SNN+Transformer \sSevenSnnTf) used different learning rates
(\sSevenLrOther vs \sSevenLrTf) and cannot support a claim.

\conf For the prospectively specified benchmark we probed S11 and S13, where the cue is further from the junction than
the view reaches. The memoryless CNN was at chance, as required (\probeElevenOne and \probeThirteenOne success), but so
was a CNN with a \memWindow-frame stack (\probeElevenTwelve and \probeThirteenTwelve), with episodes of about
\probeElevenTwelveLen--\probeThirteenTwelveLen steps (Figure~\ref{fig:probe}). Neither map met the prospectively
specified criterion, so, by the rule written before the probes, the Transformer comparison was not run. \expl A run
added after the probes, of the stacked CNN on S11 for \probeLongFrames frames, stayed at chance
(\probeElevenTwelveLong; training return \probeLongRange throughout).
\interp The episode lengths show that no agent learned to go back and look at the cue; they walk to the junction and
guess. The obstacle is exploration and credit assignment, not context length. The question whether Transformer working
memory helps is therefore inconclusive, not answered in either direction (Section~\ref{sec:limitations}).

% =====================================================================================================
\section{What did not survive controls}
\label{sec:didnt-survive}

Several results looked positive before the controls were in place. Table~\ref{tab:survive} lists them with the control
that removed them. All ``before'' numbers are recomputed from the same frozen runs, restricted to the seeds that
existed at the time; all of them are exploratory. The last row comes from the post-freeze follow-up.

\begin{table}[t]
\centering
\caption{Earlier positive findings that shrank or disappeared once a control was added. \emph{Evidence then}: the
comparison as it stood, recomputed from the frozen runs restricted to the seeds available at the time (all exploratory).
\emph{After the control}: the comparison with the control in place; for the last row, the post-freeze follow-up, which
used the 5-task CNN sequence with LwF and replay as the shared-trunk methods (Section~\ref{sec:res-shared}).
\emph{Label}: evidential status of the ``after'' result (conf.\ = confirmatory, desc.\ = descriptive). $p$: Welch
unless noted. All arms of each comparison share their seed set (Section~\ref{sec:stats}).}
\label{tab:survive}
\footnotesize
\setlength{\tabcolsep}{4pt}
\begin{tabular}{>{\raggedright\arraybackslash}p{2.7cm}>{\raggedright\arraybackslash}p{3.3cm}>{\raggedright\arraybackslash}p{4.6cm}>{\raggedright\arraybackslash}p{2.2cm}>{\raggedright\arraybackslash}p{1.6cm}}
\toprule
Earlier finding & Evidence then & After the control & Cause & Label \\
\midrule
Sleep beats replay (SNN+Transformer) & \hybSleepReplayNThreeSeeds: \hybSleepReplayNThreeCi, $p=\hybSleepReplayNThreeP$ &
seeds 1--\nHybSleep: $p=\hybSleepReplayNSixP$; budget-matched, \matchN seeds: \matchAccCi, $p=\matchAccP$ & seed count; replay-sample budget & conf. \\\addlinespace[3pt]
Shared trunk far more efficient than isolation & ACC per parameter: \perParamRatioRange{}$\times$ &
per total memory at the default buffer, isolation is \perMbRatioRange{}$\times$ more efficient; with 200 states/task
(SNN), similar mean ACC (\smallSnnTwoHundredMean vs \snnIsoMean; \smallSnnTwoHundredN seeds, wide CI) in less than half the
memory & memory accounting & desc.; conf.\ (SNN) \\\addlinespace[3pt]
Sleep far beats isolation on 5 tasks (SNN) & \framesPerTaskFlat/task, \fiveSnnShortN seeds: \fiveSnnShortSleep vs \fiveSnnShortIso, $p=\fiveSnnShortP$ &
\framesPerTaskFiveFair/task, \fiveSnnFairN seeds: \fiveSnnFairSleep vs \fiveSnnFairIso, \fiveSnnFairCi, $p=\fiveSnnFairP$ & training budget & conf. \\\addlinespace[3pt]
Homeostasis helps backprop & \homeoEarlySeeds{} (homeostasis vs backprop): \homeoEarlySolved solved, $p=\homeoEarlyP$ &
\rqOneBpN seeds: \rqOneContrCAucCi, $p=\rqOneContrCAucP$, Holm \rqOneContrCAucHolm & seed count; multiple comparisons & conf. \\\addlinespace[3pt]
Homeostasis makes DFA work & \dfaEarlySolvedHomeo vs \dfaEarlySolvedDfa solved, AUC $p$\dfaEarlyAucP &
\rqOneDfaN seeds: \rqOneDfaHomeoSolved vs \rqOneDfaSolved, Fisher $p=\dfaFisherP$ (AUC effect survives) & seed count & conf. \\\addlinespace[3pt]
Transformer helps on a memory task & S7 pilot: \sSevenSnnTf vs \sSevenSnnEight (frame stack) &
memoryless CNN on S7: \sSevenCnnOne; on S11/S13 no memory agent learns & benchmark validity; learning-rate confound & conf.\ (rule) \\\addlinespace[3pt]
With a small buffer, a shared trunk is more memory-efficient than isolation & against full-width isolation: CNN, sleep with
200 states/task \smallCnnTwoHundredPerMb vs \memCnnThreeIsoPerMb ACC per MB; SNN as in row 2 &
quarter-width isolation, CNN 5 tasks, \postN seeds: \postNarrowMean at \postNarrowMb~MB vs shared trunk with LwF
\postLwfMean at \postLwfMb~MB (\postHEightCi) and with replay \postReplayMean at \postReplayMb~MB (\postHNineCi) &
network width not matched & conf., post-freeze \\
\bottomrule
\end{tabular}
\end{table}

\paragraph{Seed count.} \interp A handful of seeds was enough to produce $p=\hybSleepReplayNThreeP$ (sleep vs replay)
and $p=\homeoEarlyP$ (homeostasis vs backprop), and a \dfaEarlySolvedHomeo vs \dfaEarlySolvedDfa split that looked
decisive (DFA). All three shrank once the seed count was raised. With bimodal outcomes, one failed seed moves a
small-sample mean a lot.

\paragraph{Training budget.} \interp At \framesPerTaskFlat frames per task, a freshly initialised SNN often had not
learned its task yet, so isolation (which starts from scratch on every task) was under-trained while the shared trunk
benefited from transfer. At a budget where fresh networks learn (\framesPerTaskFiveFair, chosen by single-task probes
before the runs), the gap shrank from \fiveSnnShortCi to \fiveSnnFairCi. Sleep still learned more tasks to
\learnThreshold or above (\fiveSnnFairLearnedSleep vs \fiveSnnFairLearnedIso), which is consistent with some positive
transfer.

\paragraph{Memory accounting.} \interp Comparing per parameter ignores the replay buffer, which is the price the
shared-trunk methods pay for not forgetting. With the buffer counted, the conclusion reversed at the default buffer
size.

\paragraph{Network width.} \interp Comparing per unit of memory still fixed isolation's width at the shared trunk's.
With quarter-width isolated networks, isolation beat both shared-trunk methods at equal or lower memory on the 5-task
CNN sequence, and differed from full width by only \postHTenCi. The small-buffer exception therefore depends on a width choice
that the frozen study did not vary.

\paragraph{Replay-sample budget.} \interp Sleep and replay are usually compared at equal buffer size, but the number of
replayed samples they consume can differ a lot (here by \unmatchedExtraPct\% on average). Matching it removed the
difference.

\paragraph{Benchmark validity.} \interp A memory benchmark must fail without memory and be learnable with it. S7 failed
the first condition; S11 and S13 failed the second at our budget.

% =====================================================================================================
\section{Limitations}
\label{sec:limitations}

\paragraph{Task identity is given.} Our setting is task-incremental: the agent knows which task it is doing. Removing
the task identity from fetch3 would make it unsolvable (the goal is not visible), so a task-agnostic study needs a
different benchmark, and sleep's distillation, which relies on per-task heads, would need redesigning.

\paragraph{Scale.} All agents are small (\paramsCnnPerTask--\paramsHybridShared parameters), the environments are
MiniGrid gridworlds, sequences have 3 or 5 tasks, and budgets are \framesPerTaskFlat--\framesPerTaskFiveFair frames per
task. Results may not transfer to larger agents, longer sequences or visually rich tasks.

\paragraph{CPU only, few seeds.} Everything ran on a 4-core CPU. Many continual arms have only \nCnnSleep seeds, so their
CIs are wide and a tie means at most ``no large difference'' where the interval is narrow, never equivalence. The
prospectively specified comparisons used \matchN--\rqOneBpN seeds per arm, and the post-freeze follow-up \postN.

\paragraph{Energy is an estimated arithmetic energy proxy.} The energy numbers multiply operation counts by literature
constants for 45\,nm arithmetic and assume event-driven hardware. We measured no power. The only cost we measured is
CPU time in our implementation, and there the SNN was slower.

\paragraph{Single settings.} Hyperparameters (homeostasis target, $\alpha$, EWC $\lambda$ beyond three values, sleep
schedule) were not tuned per arm. A different homeostasis target or sleep schedule could change the results.

\paragraph{Baselines are simple.} In the frozen study, isolation uses one full-width network per task; the post-freeze
follow-up shows that narrower isolated networks change the memory comparison (Section~\ref{sec:res-shared}), and
PackNet-style masks would be cheaper still. The shared trunk was never shrunk in the same way. Our sleep phase is replay
plus distillation, not the sleep replay consolidation of \citet{tadros2022sleep}.

\paragraph{The Transformer question is inconclusive.} On S7 memory was unnecessary; on S11 and S13 no agent, including
a CNN with a \memWindow-frame stack, learned to use the cue, even at \probeLongFrames frames. Benchmark validity was
therefore not established, and the comparison was not run.

\paragraph{Prospective specification was informal.} The test plans were recorded in the project's version-controlled
experiment log before the runs, not in a public registry. The first two phases were exploratory, and their positive
findings are exactly the ones Section~\ref{sec:didnt-survive} revisits. The follow-up was planned after the frozen
results were known, although each of its tests was specified before its own runs.

% =====================================================================================================
\ifanonymous
\section{Future work}
\label{sec:future}
These are plans and open questions, not results. The main limitation of this study is scale: every experiment ran on
a 4-core CPU, where CPU-bound environment stepping, not model size, was the bottleneck. A natural next step is to port
the agent, the mechanisms and the baselines to JAX and run them on GPUs with XLand-MiniGrid \citep{nikulin2024xland}, a
GPU-native implementation of MiniGrid with a rule-and-goal system for generating many related tasks, after first
replicating the headline results on its MiniGrid ports (sleep versus replay at a matched replay budget, STDP versus a
same-size random update, and the isolation-width result). Open questions for such a study are: longer sequences of
10--50 related tasks where transfer matters, with full accuracy-versus-memory curves in which every method shrinks along
its own size setting (isolation width, shared-trunk width, buffer size); 10--20 seeds per arm, with a power analysis and
equivalence tests, so that ``no large effect'' nulls become informative; a memory benchmark that agents can learn,
followed by Transformer versus GRU versus frame stack and a stateful SNN; the SNN version of the isolation-width
comparison at an adequate training budget (\framesPerTaskFiveFair frames per task); whether the conclusions change with
network size; and a task-agnostic setting with observable goals. Smaller follow-ups are a replay-budget sweep for sleep
versus interleaved replay (motivated by an exploratory observation in the follow-up log that sleep forgot more when
replay was scarce), separating offline replay from self-distillation within sleep, plots of the STDP runaway mechanism,
a stronger CNN sparsifier such as k-winners-take-all, compressed or generative replay \citep{shin2017continual} to
shrink the buffer, a PackNet variant that can regrow capacity, and one standard benchmark such as Continual World
\citep{wolczyk2021continual} for comparability. The methodology would stay the same: prospectively specified tests,
fresh seeds for confirmation, matched resources, and a log of every failure.
\else
\section{Future Work and Proposed Next Phase}
\label{sec:future}
This section describes plans and open questions, not results; we make no prediction about what the scaled study will
find.

\paragraph{Motivation.} The main limitation of this study is scale. Every experiment ran on a 4-core CPU, and the
bottleneck was CPU-bound environment stepping, not model size: the agents are small, and most of each run's time went
to simulating MiniGrid. This is what limited the seed counts, the sequence lengths and the training budgets.

\paragraph{Scaling plan.} Over the coming months we plan to port the agent, the mechanisms and the baselines to JAX
and run them on cloud GPU resources with XLand-MiniGrid \citep{nikulin2024xland}. \lit XLand-MiniGrid is a GPU-native
implementation of MiniGrid, reported to simulate millions of environment steps per second on one GPU, with a
rule-and-goal system for generating many related tasks.

\paragraph{Validation first.} Before trusting any new result, we will replicate the headline results on
XLand-MiniGrid's ports of the MiniGrid environments: sleep versus replay at a matched replay budget, STDP versus a
random update of the same size, and the post-freeze isolation-width result (C4). Any result that does not replicate will
be investigated before new comparisons are run.

\paragraph{Questions the scaled study will address.}
\begin{itemize}
  \item Longer task sequences (10--50 related tasks) where transfer matters, with full accuracy-versus-memory curves in
        which every method shrinks along its own size setting: isolation width, shared-trunk width and buffer size.
  \item 10--20 seeds per arm, with a power analysis before the runs and equivalence tests with stated margins, so that
        the current ``no large effect'' nulls become informative.
  \item A memory benchmark that agents can actually learn; then Transformer versus GRU versus frame stack, and a
        stateful SNN that carries its membrane state across frames.
  \item The SNN version of C4 at an adequate training budget (\framesPerTaskFiveFair frames per task).
  \item Whether the conclusions change with network size.
  \item A task-agnostic setting with observable goals, so that the task identity can be removed.
\end{itemize}

\paragraph{Smaller follow-ups.} A replay-budget sweep for sleep versus interleaved replay, motivated by an exploratory
result in the follow-up log (C3 in \texttt{explore/EXPLORE.md}) in which sleep forgot more than interleaved replay when
replay was scarce; we treat that result as exploratory only. Splitting sleep into offline replay alone versus offline
replay plus self-distillation. Plots of the STDP runaway mechanism (firing rates, weights and pairing statistics over
training). A stronger CNN sparsifier, such as k-winners-take-all. Compressed or generative replay
\citep{shin2017continual} to shrink the buffer. A PackNet variant \citep{mallya2018packnet} that can regrow capacity.
One standard benchmark, such as Continual World \citep{wolczyk2021continual}, for comparability with other work.

\paragraph{Methodology.} The same methodology throughout: tests prospectively specified in a version-controlled
experiment log, fresh seeds for confirmation, matched resources (frames, seeds, replayed samples and total memory), and
every failure logged.

\paragraph{Timeline.} Porting and replication, about 1--2 weeks; the scaled study, about 2--4 weeks; after that,
revision of this paper and journal submission, with TMLR as the first choice.
\fi

% =====================================================================================================
\section{Conclusion}
\label{sec:conclusion}

We isolated six brain-inspired mechanisms in a small task-incremental continual RL setting and gave each a matched
control; a seventh, Transformer working memory, remained inconclusive because benchmark validity was not established.
\ours No mechanism produced a statistically reliable improvement in final accuracy or forgetting beyond its matched
control; given the seed counts, these results rule out large effects rather than show equivalence, and only where the
intervals are narrow. \emph{Survived}: homeostasis sped up learning of a spiking encoder with backprop-free encoder
credit assignment (DFA; heads use exact gradients), without making it reliable. \emph{Conditional}: with a small replay
buffer, a shared trunk with per-task heads and sleep reached similar mean accuracy to full-width isolation in less than
half the memory, but in a post-freeze follow-up quarter-width isolation beat shared-trunk methods at equal or lower
memory on a 5-task CNN sequence; and the spiking encoder reached a lower operation-count regime whose energy value is
an estimated arithmetic proxy and whose measured CPU cost in our implementation was higher. \emph{Did not survive}: the
tested stabilised STDP rule against a same-size random update, homeostasis with backprop, and sleep against replay at a
matched replay budget, together with the earlier positive findings of Section~\ref{sec:didnt-survive}.
\emph{Unresolved}: the Transformer question and the SNN version of the isolation-width comparison. \interp The most
consistent lesson is methodological: in this setting, the earlier positive effects of biological mechanisms shrank or
disappeared once seed counts, training budgets, replay budgets, memory accounting and, after the freeze, network width
were made fair; the speed-up that homeostasis gives DFA survived. We suggest that claims of this kind are reported with
those budgets stated and matched.

% =====================================================================================================
\section*{Reproducibility statement}
\label{sec:reproducibility}
All code, configuration files, the frozen run directories and the analysis scripts are in \repourl. For each of the
\computeRuns PPO runs the repository keeps the full configuration, the training log, the evaluation results and the
final checkpoint, so every number can be traced to its run and every run repeated with its seed. There are three
levels of reproduction.
\begin{enumerate}
  \item \emph{The paper from the frozen results} (minutes, no training). \texttt{python paper/compute\_numbers.py}
        regenerates every number in the text; \texttt{paper/NUMBERS.md} lists each one with its source file.
        \texttt{python paper/figures/make\_figures.py} regenerates the figures, and \texttt{latexmk -pdf main.tex}
        builds the paper. The statistical tests are rerun by \texttt{scripts/rq1\_stats.py},
        \texttt{scripts/rq5\_stats.py}, \texttt{scripts/s4\_stats.py} and \texttt{scripts/memory\_fair.py}; the
        post-freeze follow-up's tests by \texttt{explore/confirm\_stats.py}, and its experiment log, with the
        prospective specification and every pilot, is \texttt{explore/EXPLORE.md}.
  \item \emph{A single run.} \texttt{python train.py -{}-config configs/<name>.yaml -{}-set seed=<s>} for a single task and
        \texttt{python continual.py -{}-suite fetch3 -{}-method <method> -{}-config configs/<name>.yaml -{}-seed <s>} for a
        task sequence; any stored configuration key can be overridden with \texttt{-{}-set}. The follow-up's runs use
        \texttt{python explore/run\_cl.py} with the same flags (the job lists are in \texttt{explore/jobs/}).
  \item \emph{All experiments.} The job lists in \texttt{jobs/} run through \texttt{scripts/run\_queue.py} (resumable,
        one thread per job). The frozen runs took about \computeCoreHours single-thread CPU hours in total, on a
        4-core machine without a GPU.
\end{enumerate}
The software environment is Python with PyTorch (CPU), Gymnasium and MiniGrid (\texttt{requirements.txt}); the unit
tests run with \texttt{pytest tests/}. The \texttt{demo/} folder holds an offline results page, GIFs of the naive and
sleep agents after the full fetch3 sequence (final accuracy-matrix rows \demoNaiveRow and \demoSleepRow), and a script
that runs a trained agent for one episode on CPU.

\section*{AI assistance disclosure}
\todo{Author to write: which parts of the code, experiments, analysis and text were produced with AI assistance, and
how they were checked.}

\bibliography{refs}
\bibliographystyle{tmlr}

% =====================================================================================================
\appendix
\section{Protocol constants}
\label{app:protocol}
\begin{center}
\small
\begin{tabular}{lp{11.5cm}}
\toprule
PPO & \numEnvs envs $\times$ \numSteps steps, \updateEpochs epochs, minibatch \minibatch, clip \clipCoef, $\gamma=\gammaVal$, $\lambda_{\text{GAE}}=\gaeLambda$, KL stop \targetKl \\
Learning rate & \lrFlat (CNN, SNN), \lrHybrid (SNN+Transformer, memory probes) \\
Frames per task & \framesPerTaskFlat (CNN, SNN), \framesPerTaskHybrid (SNN+Transformer), \framesPerTaskFiveFair (5-task SNN, fair) \\
SNN & \snnT timesteps, learnable leak (init \lifBeta), threshold 1, surrogate slope \surrogateSlope, membrane reset per frame \\
Plasticity & $\alpha=\stdpAlpha$ (prospectively specified), homeostasis target \homeoTarget, rate \homeoLr \\
Sleep & every \sleepPeriod updates + end of task, \sleepSteps steps per phase, \bufferPerTask states per task \\
Shared trunk (CNN) & trunk \paramsCnnTrunk parameters; per task one head (\paramsHead) and one task embedding (\paramsTaskEmb) \\
Post-freeze C4 & CNN, 5 tasks; quarter-width isolation: channels \postNarrowChannels, \postNarrowFeat features, \postNarrowHidden hidden units; LwF: \postLwfBatch distilled states per PPO minibatch, int8 teacher; replay: \postReplayBuf states per task \\
Evaluation & \evalEpisodesCl episodes per task (continual), \evalEpisodesSingle (single task), sampled actions \\
Memory accounting & \bytesPerState B/state (CNN, SNN), \bytesPerStateHybrid B/state (SNN+Transformer) \\
\bottomrule
\end{tabular}
\end{center}

\section{Additional numbers}
\label{app:numbers}
EWC on the CNN (3 tasks): ACC \ewcAccTen, \ewcAccHundred and \ewcAccThousand for $\lambda=10,100,1000$ (forgetting
\ewcForgetTen, \ewcForgetHundred, \ewcForgetThousand). Isolation on 5 tasks at the fair budget: \memSnnFiveFairIsoMb~MB,
ACC \memSnnFiveFairIsoAcc; sleep \memSnnFiveFairSleepMb~MB, ACC \memSnnFiveFairSleepAcc. SNN replay with small buffers:
200 states per task \smallSnnReplayTwoHundredAcc, 2000 states per task \smallSnnReplayTwoThousandAcc. The complete list
of numbers, with sources, is \texttt{paper/NUMBERS.md}.

\end{document}
